%%%%%%%%%%%%%%%%%%%%%%%%%%%%%%%%%%%%%%%%%%%%%%%%%%%%%%%%%%%%
% 8.14 BAYESIAN STATISTICS
% The Composition of Advanced Mathematics
%%%%%%%%%%%%%%%%%%%%%%%%%%%%%%%%%%%%%%%%%%%%%%%%%%%%%%%%%%%%

\section{Bayesian Statistics}
\label{sec:bayesian-statistics}

\prerequisite{
Conditional probability, Bayes' theorem, random variables, probability
distributions, expectation, variance, joint distributions, likelihood,
statistical estimation, hypothesis testing, mathematical statistics,
integration, and basic numerical methods.
}

\learningobjectives{
By the end of this section, the reader should be able to:
\begin{itemize}
    \item explain the Bayesian interpretation of probability;
    \item distinguish prior, likelihood, and posterior distributions;
    \item apply Bayes' theorem to statistical parameters;
    \item derive posterior distributions for conjugate models;
    \item compute posterior means, medians, modes, and variances;
    \item construct Bayesian credible intervals;
    \item distinguish credible intervals from frequentist confidence
          intervals;
    \item derive posterior predictive distributions;
    \item understand conjugate prior families;
    \item work with Beta--Bernoulli and Beta--Binomial models;
    \item work with Gamma--Poisson models;
    \item work with normal conjugate models;
    \item understand informative, weakly informative, and diffuse priors;
    \item understand improper priors and posterior propriety;
    \item perform prior predictive reasoning;
    \item understand Bayesian decision theory;
    \item derive Bayes estimators under common loss functions;
    \item understand Bayes factors;
    \item distinguish posterior odds from Bayes factors;
    \item understand marginal likelihoods;
    \item understand Bayesian model comparison conceptually;
    \item understand hierarchical Bayesian models;
    \item understand partial pooling;
    \item recognize empirical Bayes methods;
    \item understand exchangeability conceptually;
    \item understand Bayesian regression conceptually;
    \item understand Bayesian generalized linear models conceptually;
    \item recognize Markov chain Monte Carlo methods;
    \item understand Monte Carlo integration;
    \item recognize Gibbs sampling and Metropolis--Hastings;
    \item understand convergence diagnostics conceptually;
    \item understand posterior sensitivity to prior assumptions;
    \item connect Bayesian statistics with decision theory, regression,
          statistical learning, and computational statistics.
\end{itemize}
}

%%%%%%%%%%%%%%%%%%%%%%%%%%%%%%%%%%%%%%%%%%%%%%%%%%%%%%%%%%%%
\subsection{What Is Bayesian Statistics?}
\label{subsec:what-is-bayesian-statistics}
%%%%%%%%%%%%%%%%%%%%%%%%%%%%%%%%%%%%%%%%%%%%%%%%%%%%%%%%%%%%

Bayesian statistics treats uncertainty about unknown quantities using
probability distributions.

Suppose data are denoted

\[
x
\]

and an unknown parameter is

\[
\theta.
\]

Bayesian inference combines:

\[
\boxed{
\text{prior information}
}
\]

with

\[
\boxed{
\text{observed data}
}
\]

to produce

\[
\boxed{
\text{posterior information}.
}
\]

The basic relationship is

\[
\boxed{
\text{Posterior}
\propto
\text{Likelihood}
\times
\text{Prior}.
}
\]

%%%%%%%%%%%%%%%%%%%%%%%%%%%%%%%%%%%%%%%%%%%%%%%%%%%%%%%%%%%%
\subsection{Bayes' Theorem for Parameters}
\label{subsec:bayes-theorem-parameters}
%%%%%%%%%%%%%%%%%%%%%%%%%%%%%%%%%%%%%%%%%%%%%%%%%%%%%%%%%%%%

Let

\[
\pi(\theta)
\]

denote the prior density of \(\theta\).

Let

\[
f(x\mid\theta)
\]

denote the sampling distribution or likelihood contribution of the data.

Then the posterior density is

\[
\boxed{
\pi(\theta\mid x)
=
\frac{
f(x\mid\theta)\pi(\theta)
}{
\int
f(x\mid u)\pi(u)\,du
}.
}
\]

The denominator normalizes the posterior so that it integrates to one.

%%%%%%%%%%%%%%%%%%%%%%%%%%%%%%%%%%%%%%%%%%%%%%%%%%%%%%%%%%%%
\subsection{Posterior Proportionality}
\label{subsec:posterior-proportionality}
%%%%%%%%%%%%%%%%%%%%%%%%%%%%%%%%%%%%%%%%%%%%%%%%%%%%%%%%%%%%

Because the denominator does not depend on \(\theta\),

\[
\boxed{
\pi(\theta\mid x)
\propto
f(x\mid\theta)\pi(\theta).
}
\]

This proportional form is often the easiest starting point for
derivations.

One identifies the posterior family by collecting all terms that depend
on \(\theta\).

%%%%%%%%%%%%%%%%%%%%%%%%%%%%%%%%%%%%%%%%%%%%%%%%%%%%%%%%%%%%
\subsection{Prior Distribution}
\label{subsec:prior-distribution}
%%%%%%%%%%%%%%%%%%%%%%%%%%%%%%%%%%%%%%%%%%%%%%%%%%%%%%%%%%%%

\begin{definitionbox}{Prior Distribution}
A prior distribution describes uncertainty about an unknown parameter
before the current data are incorporated.
\end{definitionbox}

The prior may encode:

\[
\boxed{
\text{previous data},
}
\]

\[
\boxed{
\text{scientific knowledge},
}
\]

\[
\boxed{
\text{physical constraints},
}
\]

or

\[
\boxed{
\text{regularization assumptions}.
}
\]

%%%%%%%%%%%%%%%%%%%%%%%%%%%%%%%%%%%%%%%%%%%%%%%%%%%%%%%%%%%%
\subsection{Likelihood}
\label{subsec:bayesian-likelihood}
%%%%%%%%%%%%%%%%%%%%%%%%%%%%%%%%%%%%%%%%%%%%%%%%%%%%%%%%%%%%

The likelihood is

\[
\boxed{
L(\theta;x)
=
f(x\mid\theta),
}
\]

viewed as a function of \(\theta\) with observed data \(x\) fixed.

The likelihood describes how compatible different parameter values are
with the observed data.

It is not itself a probability distribution over \(\theta\).

%%%%%%%%%%%%%%%%%%%%%%%%%%%%%%%%%%%%%%%%%%%%%%%%%%%%%%%%%%%%
\subsection{Posterior Distribution}
\label{subsec:posterior-distribution}
%%%%%%%%%%%%%%%%%%%%%%%%%%%%%%%%%%%%%%%%%%%%%%%%%%%%%%%%%%%%

\begin{definitionbox}{Posterior Distribution}
The posterior distribution is the conditional distribution of the
unknown parameter after observing the data.
\end{definitionbox}

It is written

\[
\boxed{
\pi(\theta\mid x).
}
\]

All Bayesian parameter inference is obtained from this distribution.

%%%%%%%%%%%%%%%%%%%%%%%%%%%%%%%%%%%%%%%%%%%%%%%%%%%%%%%%%%%%
\subsection{Marginal Likelihood}
\label{subsec:marginal-likelihood}
%%%%%%%%%%%%%%%%%%%%%%%%%%%%%%%%%%%%%%%%%%%%%%%%%%%%%%%%%%%%

The normalizing constant

\[
\boxed{
m(x)
=
\int
f(x\mid\theta)\pi(\theta)\,d\theta
}
\]

is called the marginal likelihood or evidence.

Thus

\[
\boxed{
\pi(\theta\mid x)
=
\frac{
f(x\mid\theta)\pi(\theta)
}{
m(x)
}.
}
\]

The marginal likelihood is central to Bayesian model comparison.

%%%%%%%%%%%%%%%%%%%%%%%%%%%%%%%%%%%%%%%%%%%%%%%%%%%%%%%%%%%%
\subsection{Discrete Parameter Version}
\label{subsec:discrete-parameter-bayes}
%%%%%%%%%%%%%%%%%%%%%%%%%%%%%%%%%%%%%%%%%%%%%%%%%%%%%%%%%%%%

If \(\theta\) takes discrete values,

\[
\theta_1,\ldots,\theta_k,
\]

then

\[
\boxed{
P(\theta_j\mid x)
=
\frac{
P(x\mid\theta_j)P(\theta_j)
}{
\sum_{r=1}^{k}
P(x\mid\theta_r)P(\theta_r)
}.
}
\]

This is ordinary Bayes' theorem applied to statistical hypotheses.

%%%%%%%%%%%%%%%%%%%%%%%%%%%%%%%%%%%%%%%%%%%%%%%%%%%%%%%%%%%%
\subsection{Worked Example: Discrete Bayesian Updating}
\label{subsec:worked-discrete-bayes}
%%%%%%%%%%%%%%%%%%%%%%%%%%%%%%%%%%%%%%%%%%%%%%%%%%%%%%%%%%%%

\begin{workedexamplebox}{Two Possible Models}

Suppose

\[
P(H_1)=0.6,
\qquad
P(H_2)=0.4.
\]

Observed data \(D\) satisfy

\[
P(D\mid H_1)=0.2,
\qquad
P(D\mid H_2)=0.5.
\]

Then

\[
P(H_1\mid D)
=
\frac{
0.2(0.6)
}{
0.2(0.6)+0.5(0.4)
}.
\]

Therefore,

\begin{answerbox}
\[
\boxed{
P(H_1\mid D)
=
\frac{0.12}{0.32}
=
0.375.
}
\]

Similarly,

\[
\boxed{
P(H_2\mid D)
=
0.625.
}
\]
\end{answerbox}

\end{workedexamplebox}

%%%%%%%%%%%%%%%%%%%%%%%%%%%%%%%%%%%%%%%%%%%%%%%%%%%%%%%%%%%%
\subsection{Conjugate Priors}
\label{subsec:conjugate-priors}
%%%%%%%%%%%%%%%%%%%%%%%%%%%%%%%%%%%%%%%%%%%%%%%%%%%%%%%%%%%%

\begin{definitionbox}{Conjugate Prior}
A prior family is conjugate to a likelihood family if the posterior
belongs to the same distributional family as the prior.
\end{definitionbox}

Conjugacy produces closed-form posterior updates.

Common examples include:

\[
\boxed{
\text{Beta prior + Bernoulli likelihood},
}
\]

\[
\boxed{
\text{Gamma prior + Poisson likelihood},
}
\]

and

\[
\boxed{
\text{Normal prior + Normal likelihood}.
}
\]

%%%%%%%%%%%%%%%%%%%%%%%%%%%%%%%%%%%%%%%%%%%%%%%%%%%%%%%%%%%%
\subsection{Beta Distribution as a Prior}
\label{subsec:beta-prior}
%%%%%%%%%%%%%%%%%%%%%%%%%%%%%%%%%%%%%%%%%%%%%%%%%%%%%%%%%%%%

For a probability parameter

\[
0<p<1,
\]

a natural prior is

\[
\boxed{
p
\sim
\operatorname{Beta}(\alpha,\beta).
}
\]

Its density is

\[
\boxed{
\pi(p)
=
\frac1{B(\alpha,\beta)}
p^{\alpha-1}
(1-p)^{\beta-1}.
}
\]

Its mean is

\[
\boxed{
\mathbb{E}[p]
=
\frac{\alpha}{\alpha+\beta}.
}
\]

%%%%%%%%%%%%%%%%%%%%%%%%%%%%%%%%%%%%%%%%%%%%%%%%%%%%%%%%%%%%
\subsection{Beta--Bernoulli Model}
\label{subsec:beta-bernoulli}
%%%%%%%%%%%%%%%%%%%%%%%%%%%%%%%%%%%%%%%%%%%%%%%%%%%%%%%%%%%%

Suppose

\[
X_i\mid p
\overset{\text{i.i.d.}}{\sim}
\operatorname{Bernoulli}(p),
\]

and

\[
p
\sim
\operatorname{Beta}(\alpha,\beta).
\]

Let

\[
S
=
\sum_{i=1}^{n}X_i.
\]

The likelihood is

\[
\boxed{
L(p)
\propto
p^S
(1-p)^{n-S}.
}
\]

%%%%%%%%%%%%%%%%%%%%%%%%%%%%%%%%%%%%%%%%%%%%%%%%%%%%%%%%%%%%
\subsection{Deriving the Beta Posterior}
\label{subsec:derive-beta-posterior}
%%%%%%%%%%%%%%%%%%%%%%%%%%%%%%%%%%%%%%%%%%%%%%%%%%%%%%%%%%%%

Multiply prior and likelihood:

\[
\pi(p\mid x)
\propto
p^S
(1-p)^{n-S}
p^{\alpha-1}
(1-p)^{\beta-1}.
\]

Therefore,

\[
\pi(p\mid x)
\propto
p^{\alpha+S-1}
(1-p)^{\beta+n-S-1}.
\]

Hence,

\[
\boxed{
p\mid x
\sim
\operatorname{Beta}
(
\alpha+S,
\beta+n-S
).
}
\]

%%%%%%%%%%%%%%%%%%%%%%%%%%%%%%%%%%%%%%%%%%%%%%%%%%%%%%%%%%%%
\subsection{Beta Posterior Mean}
\label{subsec:beta-posterior-mean}
%%%%%%%%%%%%%%%%%%%%%%%%%%%%%%%%%%%%%%%%%%%%%%%%%%%%%%%%%%%%

The posterior mean is

\[
\boxed{
\mathbb{E}[p\mid x]
=
\frac{
\alpha+S
}{
\alpha+\beta+n
}.
}
\]

This can be rewritten as a weighted average:

\[
\boxed{
\mathbb{E}[p\mid x]
=
\frac{
\alpha+\beta
}{
\alpha+\beta+n
}
\frac{\alpha}{\alpha+\beta}
+
\frac{
n
}{
\alpha+\beta+n
}
\frac{S}{n}.
}
\]

Thus the posterior mean combines the prior mean and sample proportion.

%%%%%%%%%%%%%%%%%%%%%%%%%%%%%%%%%%%%%%%%%%%%%%%%%%%%%%%%%%%%
\subsection{Prior Effective Sample Size}
\label{subsec:beta-effective-sample-size}
%%%%%%%%%%%%%%%%%%%%%%%%%%%%%%%%%%%%%%%%%%%%%%%%%%%%%%%%%%%%

In the Beta--Bernoulli model, the quantity

\[
\boxed{
\alpha+\beta
}
\]

is often interpreted as a prior strength or prior effective sample size.

This interpretation is useful but depends on the model and
parameterization.

Larger values make the prior more concentrated.

%%%%%%%%%%%%%%%%%%%%%%%%%%%%%%%%%%%%%%%%%%%%%%%%%%%%%%%%%%%%
\subsection{Worked Example: Beta--Bernoulli Updating}
\label{subsec:worked-beta-bernoulli}
%%%%%%%%%%%%%%%%%%%%%%%%%%%%%%%%%%%%%%%%%%%%%%%%%%%%%%%%%%%%

\begin{workedexamplebox}{Updating a Probability}

Suppose

\[
p
\sim
\operatorname{Beta}(2,2).
\]

Observe

\[
n=10
\]

trials with

\[
S=7
\]

successes.

Then

\[
p\mid x
\sim
\operatorname{Beta}
(
2+7,
2+3
).
\]

Therefore,

\begin{answerbox}
\[
\boxed{
p\mid x
\sim
\operatorname{Beta}(9,5).
}
\]

The posterior mean is

\[
\boxed{
\frac9{14}
\approx
0.643.
}
\]
\end{answerbox}

\end{workedexamplebox}

%%%%%%%%%%%%%%%%%%%%%%%%%%%%%%%%%%%%%%%%%%%%%%%%%%%%%%%%%%%%
\subsection{Posterior Mode}
\label{subsec:posterior-mode}
%%%%%%%%%%%%%%%%%%%%%%%%%%%%%%%%%%%%%%%%%%%%%%%%%%%%%%%%%%%%

The posterior mode is also called the maximum a posteriori estimate,
abbreviated MAP.

It is

\[
\boxed{
\widehat{\theta}_{MAP}
=
\arg\max_\theta
\pi(\theta\mid x).
}
\]

Because

\[
\pi(\theta\mid x)
\propto
L(\theta)\pi(\theta),
\]

the MAP incorporates both likelihood and prior information.

%%%%%%%%%%%%%%%%%%%%%%%%%%%%%%%%%%%%%%%%%%%%%%%%%%%%%%%%%%%%
\subsection{MLE Versus MAP}
\label{subsec:mle-versus-map}
%%%%%%%%%%%%%%%%%%%%%%%%%%%%%%%%%%%%%%%%%%%%%%%%%%%%%%%%%%%%

The MLE maximizes

\[
\boxed{
L(\theta).
}
\]

The MAP maximizes

\[
\boxed{
L(\theta)\pi(\theta).
}
\]

Therefore:

\[
\boxed{
\text{MLE}
=
\text{data-only likelihood maximizer},
}
\]

while

\[
\boxed{
\text{MAP}
=
\text{posterior maximizer}.
}
\]

With a flat prior over the relevant region, the two may coincide.

%%%%%%%%%%%%%%%%%%%%%%%%%%%%%%%%%%%%%%%%%%%%%%%%%%%%%%%%%%%%
\subsection{Posterior Median}
\label{subsec:posterior-median}
%%%%%%%%%%%%%%%%%%%%%%%%%%%%%%%%%%%%%%%%%%%%%%%%%%%%%%%%%%%%

The posterior median \(m\) satisfies

\[
\boxed{
P(\theta\leq m\mid x)
\geq
\frac12
}
\]

and

\[
\boxed{
P(\theta\geq m\mid x)
\geq
\frac12.
}
\]

It can differ from both the posterior mean and posterior mode when the
posterior is skewed.

%%%%%%%%%%%%%%%%%%%%%%%%%%%%%%%%%%%%%%%%%%%%%%%%%%%%%%%%%%%%
\subsection{Bayes Estimators and Loss Functions}
\label{subsec:bayes-estimators-loss}
%%%%%%%%%%%%%%%%%%%%%%%%%%%%%%%%%%%%%%%%%%%%%%%%%%%%%%%%%%%%

A Bayesian point estimate depends on the chosen loss function.

Under squared-error loss,

\[
\boxed{
L(\theta,a)
=
(\theta-a)^2,
}
\]

the Bayes estimator is the posterior mean.

Under absolute-error loss,

\[
\boxed{
L(\theta,a)
=
|\theta-a|,
}
\]

a posterior median is optimal.

Under zero-one-type loss, the posterior mode arises naturally.

%%%%%%%%%%%%%%%%%%%%%%%%%%%%%%%%%%%%%%%%%%%%%%%%%%%%%%%%%%%%
\subsection{Credible Intervals}
\label{subsec:credible-intervals}
%%%%%%%%%%%%%%%%%%%%%%%%%%%%%%%%%%%%%%%%%%%%%%%%%%%%%%%%%%%%

A Bayesian \(100(1-\alpha)\%\) credible interval is an interval

\[
[C_L,C_U]
\]

satisfying

\[
\boxed{
P
(
C_L\leq\theta\leq C_U
\mid x
)
=
1-\alpha.
}
\]

This is a direct posterior probability statement about the parameter.

%%%%%%%%%%%%%%%%%%%%%%%%%%%%%%%%%%%%%%%%%%%%%%%%%%%%%%%%%%%%
\subsection{Equal-Tailed Credible Interval}
\label{subsec:equal-tailed-credible-interval}
%%%%%%%%%%%%%%%%%%%%%%%%%%%%%%%%%%%%%%%%%%%%%%%%%%%%%%%%%%%%

An equal-tailed interval satisfies

\[
\boxed{
P
(
\theta<C_L
\mid x
)
=
\frac{\alpha}{2},
}
\]

and

\[
\boxed{
P
(
\theta>C_U
\mid x
)
=
\frac{\alpha}{2}.
}
\]

The endpoints are posterior quantiles.

%%%%%%%%%%%%%%%%%%%%%%%%%%%%%%%%%%%%%%%%%%%%%%%%%%%%%%%%%%%%
\subsection{Highest Posterior Density Region}
\label{subsec:highest-posterior-density}
%%%%%%%%%%%%%%%%%%%%%%%%%%%%%%%%%%%%%%%%%%%%%%%%%%%%%%%%%%%%

A highest posterior density region, abbreviated HPD, contains parameter
values with the greatest posterior density until the desired posterior
probability is reached.

For a continuous unimodal posterior, an HPD interval satisfies
conceptually

\[
\boxed{
\pi(\theta\mid x)
\geq
c
}
\]

for all points in the region, with \(c\) chosen so that the region has
posterior probability \(1-\alpha\).

%%%%%%%%%%%%%%%%%%%%%%%%%%%%%%%%%%%%%%%%%%%%%%%%%%%%%%%%%%%%
\subsection{Credible Versus Confidence Intervals}
\label{subsec:credible-versus-confidence}
%%%%%%%%%%%%%%%%%%%%%%%%%%%%%%%%%%%%%%%%%%%%%%%%%%%%%%%%%%%%

A Bayesian credible interval supports a statement such as

\[
\boxed{
P
(
\theta\in[C_L,C_U]
\mid x
)
=
0.95.
}
\]

A frequentist confidence interval is interpreted through repeated-sample
coverage:

\[
\boxed{
P_\theta
(
\theta\in C(X)
)
=
0.95.
}
\]

These interpretations are not the same.

%%%%%%%%%%%%%%%%%%%%%%%%%%%%%%%%%%%%%%%%%%%%%%%%%%%%%%%%%%%%
\subsection{Posterior Predictive Distribution}
\label{subsec:posterior-predictive}
%%%%%%%%%%%%%%%%%%%%%%%%%%%%%%%%%%%%%%%%%%%%%%%%%%%%%%%%%%%%

Suppose a future observation is

\[
X_{\text{new}}.
\]

The posterior predictive distribution integrates over posterior
uncertainty in \(\theta\):

\[
\boxed{
p
(
x_{\text{new}}
\mid x
)
=
\int
p
(
x_{\text{new}}
\mid\theta
)
\pi(\theta\mid x)
\,d\theta.
}
\]

Thus prediction averages over possible parameter values rather than
plugging in only one estimate.

%%%%%%%%%%%%%%%%%%%%%%%%%%%%%%%%%%%%%%%%%%%%%%%%%%%%%%%%%%%%
\subsection{Beta--Bernoulli Predictive Probability}
\label{subsec:beta-bernoulli-predictive}
%%%%%%%%%%%%%%%%%%%%%%%%%%%%%%%%%%%%%%%%%%%%%%%%%%%%%%%%%%%%

If

\[
p\mid x
\sim
\operatorname{Beta}
(
\alpha+S,
\beta+n-S
),
\]

then

\[
P
(
X_{\text{new}}=1
\mid x
)
=
\mathbb{E}[p\mid x].
\]

Therefore,

\[
\boxed{
P
(
X_{\text{new}}=1
\mid x
)
=
\frac{
\alpha+S
}{
\alpha+\beta+n
}.
}
\]

%%%%%%%%%%%%%%%%%%%%%%%%%%%%%%%%%%%%%%%%%%%%%%%%%%%%%%%%%%%%
\subsection{Beta--Binomial Predictive Distribution}
\label{subsec:beta-binomial-predictive}
%%%%%%%%%%%%%%%%%%%%%%%%%%%%%%%%%%%%%%%%%%%%%%%%%%%%%%%%%%%%

Suppose \(M\) future Bernoulli trials are observed and

\[
Y
=
\text{number of future successes}.
\]

Integrating over the Beta posterior gives a Beta--Binomial distribution.

If the posterior parameters are \(a\) and \(b\), then

\[
\boxed{
P(Y=y\mid x)
=
\binom{M}{y}
\frac{
B(a+y,b+M-y)
}{
B(a,b)
}.
}
\]

%%%%%%%%%%%%%%%%%%%%%%%%%%%%%%%%%%%%%%%%%%%%%%%%%%%%%%%%%%%%
\subsection{Gamma Distribution as a Prior}
\label{subsec:gamma-prior}
%%%%%%%%%%%%%%%%%%%%%%%%%%%%%%%%%%%%%%%%%%%%%%%%%%%%%%%%%%%%

For a positive rate parameter

\[
\lambda>0,
\]

a common prior is

\[
\boxed{
\lambda
\sim
\operatorname{Gamma}(\alpha,\beta),
}
\]

using a shape--rate parameterization with density

\[
\boxed{
\pi(\lambda)
=
\frac{
\beta^\alpha
}{
\Gamma(\alpha)
}
\lambda^{\alpha-1}
e^{-\beta\lambda}.
}
\]

Then

\[
\boxed{
\mathbb{E}[\lambda]
=
\frac{\alpha}{\beta}.
}
\]

%%%%%%%%%%%%%%%%%%%%%%%%%%%%%%%%%%%%%%%%%%%%%%%%%%%%%%%%%%%%
\subsection{Gamma--Poisson Model}
\label{subsec:gamma-poisson-model}
%%%%%%%%%%%%%%%%%%%%%%%%%%%%%%%%%%%%%%%%%%%%%%%%%%%%%%%%%%%%

Suppose

\[
X_i\mid\lambda
\overset{\text{i.i.d.}}{\sim}
\operatorname{Poisson}(\lambda),
\]

and

\[
\lambda
\sim
\operatorname{Gamma}(\alpha,\beta).
\]

Let

\[
S
=
\sum_{i=1}^{n}X_i.
\]

The likelihood is

\[
\boxed{
L(\lambda)
\propto
\lambda^S
e^{-n\lambda}.
}
\]

%%%%%%%%%%%%%%%%%%%%%%%%%%%%%%%%%%%%%%%%%%%%%%%%%%%%%%%%%%%%
\subsection{Deriving the Gamma Posterior}
\label{subsec:derive-gamma-posterior}
%%%%%%%%%%%%%%%%%%%%%%%%%%%%%%%%%%%%%%%%%%%%%%%%%%%%%%%%%%%%

Multiply prior and likelihood:

\[
\pi(\lambda\mid x)
\propto
\lambda^S
e^{-n\lambda}
\lambda^{\alpha-1}
e^{-\beta\lambda}.
\]

Therefore,

\[
\pi(\lambda\mid x)
\propto
\lambda^{\alpha+S-1}
e^{-(\beta+n)\lambda}.
\]

Hence,

\[
\boxed{
\lambda\mid x
\sim
\operatorname{Gamma}
(
\alpha+S,
\beta+n
).
}
\]

%%%%%%%%%%%%%%%%%%%%%%%%%%%%%%%%%%%%%%%%%%%%%%%%%%%%%%%%%%%%
\subsection{Worked Example: Gamma--Poisson Updating}
\label{subsec:worked-gamma-poisson}
%%%%%%%%%%%%%%%%%%%%%%%%%%%%%%%%%%%%%%%%%%%%%%%%%%%%%%%%%%%%

\begin{workedexamplebox}{Updating a Poisson Rate}

Suppose

\[
\lambda
\sim
\operatorname{Gamma}(3,2),
\]

using shape--rate form.

Observe

\[
n=5
\]

counts with total

\[
S=12.
\]

Then

\[
\lambda\mid x
\sim
\operatorname{Gamma}
(
15,7
).
\]

Therefore,

\begin{answerbox}
\[
\boxed{
\mathbb{E}
[
\lambda\mid x
]
=
\frac{15}{7}
\approx
2.143.
}
\]
\end{answerbox}

\end{workedexamplebox}

%%%%%%%%%%%%%%%%%%%%%%%%%%%%%%%%%%%%%%%%%%%%%%%%%%%%%%%%%%%%
\subsection{Normal Prior for a Normal Mean}
\label{subsec:normal-normal-model}
%%%%%%%%%%%%%%%%%%%%%%%%%%%%%%%%%%%%%%%%%%%%%%%%%%%%%%%%%%%%

Suppose

\[
X_i\mid\mu
\overset{\text{i.i.d.}}{\sim}
N(\mu,\sigma^2),
\]

where \(\sigma^2\) is known.

Let

\[
\boxed{
\mu
\sim
N(\mu_0,\tau_0^2).
}
\]

The posterior is also normal.

%%%%%%%%%%%%%%%%%%%%%%%%%%%%%%%%%%%%%%%%%%%%%%%%%%%%%%%%%%%%
\subsection{Posterior Precision}
\label{subsec:posterior-precision-normal}
%%%%%%%%%%%%%%%%%%%%%%%%%%%%%%%%%%%%%%%%%%%%%%%%%%%%%%%%%%%%

Define precision as inverse variance.

The prior precision is

\[
\boxed{
\frac1{\tau_0^2}.
}
\]

The data precision for the sample mean is

\[
\boxed{
\frac{n}{\sigma^2}.
}
\]

The posterior precision is the sum:

\[
\boxed{
\frac1{\tau_n^2}
=
\frac1{\tau_0^2}
+
\frac{n}{\sigma^2}.
}
\]

Thus

\[
\boxed{
\tau_n^2
=
\left(
\frac1{\tau_0^2}
+
\frac{n}{\sigma^2}
\right)^{-1}.
}
\]

%%%%%%%%%%%%%%%%%%%%%%%%%%%%%%%%%%%%%%%%%%%%%%%%%%%%%%%%%%%%
\subsection{Normal Posterior Mean}
\label{subsec:normal-posterior-mean}
%%%%%%%%%%%%%%%%%%%%%%%%%%%%%%%%%%%%%%%%%%%%%%%%%%%%%%%%%%%%

The posterior mean is

\[
\boxed{
\mu_n
=
\tau_n^2
\left(
\frac{\mu_0}{\tau_0^2}
+
\frac{
n\overline{x}
}{
\sigma^2
}
\right).
}
\]

Equivalently, it is a precision-weighted average of the prior mean and
sample mean.

Thus

\[
\boxed{
\mu\mid x
\sim
N(\mu_n,\tau_n^2).
}
\]

%%%%%%%%%%%%%%%%%%%%%%%%%%%%%%%%%%%%%%%%%%%%%%%%%%%%%%%%%%%%
\subsection{Worked Example: Normal--Normal Updating}
\label{subsec:worked-normal-normal}
%%%%%%%%%%%%%%%%%%%%%%%%%%%%%%%%%%%%%%%%%%%%%%%%%%%%%%%%%%%%

\begin{workedexamplebox}{Combining Prior and Sample Information}

Suppose

\[
\mu
\sim
N(10,4),
\]

so

\[
\tau_0^2=4.
\]

Suppose

\[
\sigma^2=9,
\qquad
n=9,
\qquad
\overline{x}=13.
\]

The prior precision is

\[
\frac14.
\]

The data precision is

\[
\frac9{9}=1.
\]

Thus the posterior variance is

\[
\tau_n^2
=
\frac1{
1/4+1
}
=
\frac45.
\]

The posterior mean is

\[
\mu_n
=
\frac45
\left[
\frac{10}{4}
+
13
\right].
\]

Therefore,

\begin{answerbox}
\[
\boxed{
\mu_n
=
12.4,
\qquad
\tau_n^2
=
0.8.
}
\]
\end{answerbox}

\end{workedexamplebox}

%%%%%%%%%%%%%%%%%%%%%%%%%%%%%%%%%%%%%%%%%%%%%%%%%%%%%%%%%%%%
\subsection{Shrinkage}
\label{subsec:bayesian-shrinkage}
%%%%%%%%%%%%%%%%%%%%%%%%%%%%%%%%%%%%%%%%%%%%%%%%%%%%%%%%%%%%

The normal posterior mean lies between the prior mean and sample mean.

Thus the data estimate

\[
\overline{x}
\]

is pulled toward

\[
\mu_0.
\]

This phenomenon is called shrinkage.

The amount of shrinkage depends on relative prior and data precision.

%%%%%%%%%%%%%%%%%%%%%%%%%%%%%%%%%%%%%%%%%%%%%%%%%%%%%%%%%%%%
\subsection{Large-Sample Behavior}
\label{subsec:bayesian-large-sample}
%%%%%%%%%%%%%%%%%%%%%%%%%%%%%%%%%%%%%%%%%%%%%%%%%%%%%%%%%%%%

As sample size grows,

\[
\frac{n}{\sigma^2}
\]

often dominates the fixed prior precision.

Then the posterior becomes increasingly concentrated around values
favored by the likelihood.

Under regularity conditions,

\[
\boxed{
\text{posterior influence of a fixed regular prior decreases as }n
\text{ grows}.
}
\]

%%%%%%%%%%%%%%%%%%%%%%%%%%%%%%%%%%%%%%%%%%%%%%%%%%%%%%%%%%%%
\subsection{Informative Priors}
\label{subsec:informative-priors}
%%%%%%%%%%%%%%%%%%%%%%%%%%%%%%%%%%%%%%%%%%%%%%%%%%%%%%%%%%%%

An informative prior places substantial probability on a relatively
limited parameter region based on genuine prior knowledge.

It may represent:

\[
\boxed{
\text{previous studies},
}
\]

\[
\boxed{
\text{physical constraints},
}
\]

or

\[
\boxed{
\text{expert knowledge}.
}
\]

Informative priors can substantially affect posterior inference when data
are limited.

%%%%%%%%%%%%%%%%%%%%%%%%%%%%%%%%%%%%%%%%%%%%%%%%%%%%%%%%%%%%
\subsection{Weakly Informative Priors}
\label{subsec:weakly-informative-priors}
%%%%%%%%%%%%%%%%%%%%%%%%%%%%%%%%%%%%%%%%%%%%%%%%%%%%%%%%%%%%

A weakly informative prior rules out implausible parameter values while
remaining broad over scientifically reasonable values.

Such priors can stabilize estimation without forcing the posterior near a
very narrow prior belief.

They are often preferable to extremely diffuse priors in complex models.

%%%%%%%%%%%%%%%%%%%%%%%%%%%%%%%%%%%%%%%%%%%%%%%%%%%%%%%%%%%%
\subsection{Diffuse Priors}
\label{subsec:diffuse-priors}
%%%%%%%%%%%%%%%%%%%%%%%%%%%%%%%%%%%%%%%%%%%%%%%%%%%%%%%%%%%%

A diffuse prior spreads probability across a broad parameter region.

It is intended to contribute relatively little information.

However,

\[
\boxed{
\text{very broad}
\neq
\text{automatically noninformative}.
}
\]

The effect of a prior depends on parameterization and the model.

%%%%%%%%%%%%%%%%%%%%%%%%%%%%%%%%%%%%%%%%%%%%%%%%%%%%%%%%%%%%
\subsection{Improper Priors}
\label{subsec:improper-priors}
%%%%%%%%%%%%%%%%%%%%%%%%%%%%%%%%%%%%%%%%%%%%%%%%%%%%%%%%%%%%

An improper prior does not integrate to one.

For example,

\[
\boxed{
\pi(\mu)\propto1,
\qquad
-\infty<\mu<\infty,
}
\]

is not a proper probability density.

Improper priors can sometimes produce proper posterior distributions.

However, they require care because ordinary prior probability
interpretations and some model-comparison methods may no longer apply.

%%%%%%%%%%%%%%%%%%%%%%%%%%%%%%%%%%%%%%%%%%%%%%%%%%%%%%%%%%%%
\subsection{Posterior Propriety}
\label{subsec:posterior-propriety}
%%%%%%%%%%%%%%%%%%%%%%%%%%%%%%%%%%%%%%%%%%%%%%%%%%%%%%%%%%%%

A posterior is proper if

\[
\boxed{
\int
L(\theta)\pi(\theta)
\,d\theta
<
\infty.
}
\]

Then the posterior can be normalized into a valid probability
distribution.

When improper priors are used, posterior propriety must be verified.

%%%%%%%%%%%%%%%%%%%%%%%%%%%%%%%%%%%%%%%%%%%%%%%%%%%%%%%%%%%%
\subsection{Prior Predictive Distribution}
\label{subsec:prior-predictive}
%%%%%%%%%%%%%%%%%%%%%%%%%%%%%%%%%%%%%%%%%%%%%%%%%%%%%%%%%%%%

Before observing data, the prior predictive distribution is

\[
\boxed{
p(x)
=
\int
p(x\mid\theta)\pi(\theta)\,d\theta.
}
\]

It describes datasets implied jointly by the prior and likelihood.

Prior predictive simulation is useful for checking whether the model
generates plausible observations.

%%%%%%%%%%%%%%%%%%%%%%%%%%%%%%%%%%%%%%%%%%%%%%%%%%%%%%%%%%%%
\subsection{Prior Predictive Checks}
\label{subsec:prior-predictive-checks}
%%%%%%%%%%%%%%%%%%%%%%%%%%%%%%%%%%%%%%%%%%%%%%%%%%%%%%%%%%%%

A prior may look reasonable when viewed only on the parameter scale but
produce absurd predictions on the data scale.

Prior predictive checking asks:

\[
\boxed{
\text{What kinds of datasets does this prior-likelihood combination
generate?}
}
\]

If generated data are scientifically implausible, the prior or model may
need revision.

%%%%%%%%%%%%%%%%%%%%%%%%%%%%%%%%%%%%%%%%%%%%%%%%%%%%%%%%%%%%
\subsection{Posterior Predictive Checks}
\label{subsec:posterior-predictive-checks}
%%%%%%%%%%%%%%%%%%%%%%%%%%%%%%%%%%%%%%%%%%%%%%%%%%%%%%%%%%%%

Posterior predictive checks compare observed data with replicated data

\[
X^{rep}
\]

generated from

\[
\boxed{
p
(
x^{rep}\mid x
).
}
\]

One can compare statistics such as:

\[
\boxed{
\text{means},
}
\]

\[
\boxed{
\text{variances},
}
\]

\[
\boxed{
\text{extremes},
}
\]

or

\[
\boxed{
\text{graphical patterns}.
}
\]

Large discrepancies may reveal model inadequacy.

%%%%%%%%%%%%%%%%%%%%%%%%%%%%%%%%%%%%%%%%%%%%%%%%%%%%%%%%%%%%
\subsection{Bayesian Model Checking}
\label{subsec:bayesian-model-checking}
%%%%%%%%%%%%%%%%%%%%%%%%%%%%%%%%%%%%%%%%%%%%%%%%%%%%%%%%%%%%

Bayesian inference conditions on the assumed model.

Therefore posterior precision does not guarantee model correctness.

Model checking should investigate whether the model reproduces important
features of the observed data.

Thus:

\[
\boxed{
\text{posterior certainty}
\neq
\text{model validity}.
}
\]

%%%%%%%%%%%%%%%%%%%%%%%%%%%%%%%%%%%%%%%%%%%%%%%%%%%%%%%%%%%%
\subsection{Posterior Odds}
\label{subsec:posterior-odds}
%%%%%%%%%%%%%%%%%%%%%%%%%%%%%%%%%%%%%%%%%%%%%%%%%%%%%%%%%%%%

For two hypotheses

\[
H_1
\]

and

\[
H_0,
\]

posterior odds are

\[
\boxed{
\frac{
P(H_1\mid x)
}{
P(H_0\mid x)
}.
}
\]

Bayes' theorem gives

\[
\boxed{
\frac{
P(H_1\mid x)
}{
P(H_0\mid x)
}
=
\frac{
P(H_1)
}{
P(H_0)
}
\times
\frac{
P(x\mid H_1)
}{
P(x\mid H_0)
}.
}
\]

%%%%%%%%%%%%%%%%%%%%%%%%%%%%%%%%%%%%%%%%%%%%%%%%%%%%%%%%%%%%
\subsection{Bayes Factor}
\label{subsec:bayes-factor}
%%%%%%%%%%%%%%%%%%%%%%%%%%%%%%%%%%%%%%%%%%%%%%%%%%%%%%%%%%%%

The Bayes factor comparing \(H_1\) with \(H_0\) is

\[
\boxed{
BF_{10}
=
\frac{
P(x\mid H_1)
}{
P(x\mid H_0)
}.
}
\]

Thus

\[
\boxed{
\text{posterior odds}
=
\text{prior odds}
\times
\text{Bayes factor}.
}
\]

The Bayes factor measures how the observed data change the relative
support for the two models.

%%%%%%%%%%%%%%%%%%%%%%%%%%%%%%%%%%%%%%%%%%%%%%%%%%%%%%%%%%%%
\subsection{Marginal Likelihood Under a Model}
\label{subsec:model-marginal-likelihood}
%%%%%%%%%%%%%%%%%%%%%%%%%%%%%%%%%%%%%%%%%%%%%%%%%%%%%%%%%%%%

Suppose model \(M_k\) has parameter \(\theta_k\).

Its marginal likelihood is

\[
\boxed{
p(x\mid M_k)
=
\int
p(x\mid\theta_k,M_k)
\pi(\theta_k\mid M_k)
\,d\theta_k.
}
\]

Thus a Bayes factor compares integrated likelihoods rather than maximized
likelihoods.

%%%%%%%%%%%%%%%%%%%%%%%%%%%%%%%%%%%%%%%%%%%%%%%%%%%%%%%%%%%%
\subsection{Occam Effect in Bayesian Model Comparison}
\label{subsec:bayesian-occam-effect}
%%%%%%%%%%%%%%%%%%%%%%%%%%%%%%%%%%%%%%%%%%%%%%%%%%%%%%%%%%%%

A flexible model may achieve a large maximum likelihood over a narrow
parameter region.

But the marginal likelihood averages the likelihood across the prior
parameter space.

Thus models spreading prior mass over large regions unsupported by the
data may be penalized automatically.

This is sometimes called a Bayesian Occam effect.

%%%%%%%%%%%%%%%%%%%%%%%%%%%%%%%%%%%%%%%%%%%%%%%%%%%%%%%%%%%%
\subsection{Bayes Factors Depend on Priors}
\label{subsec:bayes-factor-prior-sensitivity}
%%%%%%%%%%%%%%%%%%%%%%%%%%%%%%%%%%%%%%%%%%%%%%%%%%%%%%%%%%%%

Bayes factors can be highly sensitive to prior distributions on
model-specific parameters.

Very diffuse priors may strongly reduce marginal likelihoods.

Therefore model-comparison priors require especially careful
justification.

Improper priors generally cannot be inserted naively into ordinary Bayes
factor calculations.

%%%%%%%%%%%%%%%%%%%%%%%%%%%%%%%%%%%%%%%%%%%%%%%%%%%%%%%%%%%%
\subsection{Bayesian Decision Theory}
\label{subsec:bayesian-decision-theory}
%%%%%%%%%%%%%%%%%%%%%%%%%%%%%%%%%%%%%%%%%%%%%%%%%%%%%%%%%%%%

Let

\[
a
\]

be an action and

\[
L(\theta,a)
\]

the loss.

After observing \(x\), the posterior expected loss is

\[
\boxed{
\rho(a\mid x)
=
\mathbb{E}
[
L(\theta,a)
\mid x
].
}
\]

A Bayes action minimizes posterior expected loss:

\[
\boxed{
a^*(x)
=
\arg\min_a
\rho(a\mid x).
}
\]

%%%%%%%%%%%%%%%%%%%%%%%%%%%%%%%%%%%%%%%%%%%%%%%%%%%%%%%%%%%%
\subsection{Posterior Mean Under Squared Error}
\label{subsec:posterior-mean-squared-error}
%%%%%%%%%%%%%%%%%%%%%%%%%%%%%%%%%%%%%%%%%%%%%%%%%%%%%%%%%%%%

For

\[
L(\theta,a)
=
(\theta-a)^2,
\]

posterior expected loss is minimized at

\[
\boxed{
a^*
=
\mathbb{E}[\theta\mid x].
}
\]

Thus the posterior mean is not merely a summary statistic; it is an
optimal decision under squared-error loss.

%%%%%%%%%%%%%%%%%%%%%%%%%%%%%%%%%%%%%%%%%%%%%%%%%%%%%%%%%%%%
\subsection{Posterior Median Under Absolute Error}
\label{subsec:posterior-median-absolute-error}
%%%%%%%%%%%%%%%%%%%%%%%%%%%%%%%%%%%%%%%%%%%%%%%%%%%%%%%%%%%%

For

\[
L(\theta,a)
=
|\theta-a|,
\]

posterior expected loss is minimized by a posterior median:

\[
\boxed{
a^*
=
\operatorname{Median}(\theta\mid x).
}
\]

This parallels the frequentist robustness of median-based estimation.

%%%%%%%%%%%%%%%%%%%%%%%%%%%%%%%%%%%%%%%%%%%%%%%%%%%%%%%%%%%%
\subsection{Asymmetric Loss}
\label{subsec:asymmetric-bayesian-loss}
%%%%%%%%%%%%%%%%%%%%%%%%%%%%%%%%%%%%%%%%%%%%%%%%%%%%%%%%%%%%

Real decisions may assign different costs to overestimation and
underestimation.

For example,

\[
\boxed{
L(\theta,a)
=
\begin{cases}
c_1(\theta-a),&a<\theta,\\
c_2(a-\theta),&a\geq\theta.
\end{cases}
}
\]

The optimal Bayes action is then a posterior quantile determined by the
relative costs.

Thus Bayesian decision theory naturally incorporates asymmetric
consequences.

%%%%%%%%%%%%%%%%%%%%%%%%%%%%%%%%%%%%%%%%%%%%%%%%%%%%%%%%%%%%
\subsection{Hierarchical Bayesian Models}
\label{subsec:hierarchical-bayesian-models}
%%%%%%%%%%%%%%%%%%%%%%%%%%%%%%%%%%%%%%%%%%%%%%%%%%%%%%%%%%%%

A hierarchical model introduces parameters that themselves depend on
higher-level parameters.

For groups

\[
j=1,\ldots,J,
\]

one might write

\[
\boxed{
Y_{ij}
\mid
\theta_j
\sim
N(\theta_j,\sigma^2),
}
\]

\[
\boxed{
\theta_j
\mid
\mu,\tau^2
\sim
N(\mu,\tau^2).
}
\]

The hyperparameters are

\[
\mu
\]

and

\[
\tau^2.
\]

%%%%%%%%%%%%%%%%%%%%%%%%%%%%%%%%%%%%%%%%%%%%%%%%%%%%%%%%%%%%
\subsection{Hyperparameters}
\label{subsec:hyperparameters}
%%%%%%%%%%%%%%%%%%%%%%%%%%%%%%%%%%%%%%%%%%%%%%%%%%%%%%%%%%%%

Parameters controlling prior distributions are called hyperparameters.

For example, in

\[
\theta_j
\sim
N(\mu,\tau^2),
\]

the quantities

\[
\boxed{
\mu
}
\]

and

\[
\boxed{
\tau^2
}
\]

are hyperparameters.

In a fully Bayesian model, they may themselves receive prior
distributions.

%%%%%%%%%%%%%%%%%%%%%%%%%%%%%%%%%%%%%%%%%%%%%%%%%%%%%%%%%%%%
\subsection{Partial Pooling}
\label{subsec:partial-pooling}
%%%%%%%%%%%%%%%%%%%%%%%%%%%%%%%%%%%%%%%%%%%%%%%%%%%%%%%%%%%%

Hierarchical models often produce partial pooling.

Group-specific estimates are pulled toward a common population mean.

Small groups are generally shrunk more strongly.

Large groups with substantial data are influenced more strongly by their
own observations.

Thus hierarchical modeling balances:

\[
\boxed{
\text{complete pooling}
}
\]

and

\[
\boxed{
\text{no pooling}.
}
\]

%%%%%%%%%%%%%%%%%%%%%%%%%%%%%%%%%%%%%%%%%%%%%%%%%%%%%%%%%%%%
\subsection{Complete Pooling}
\label{subsec:complete-pooling}
%%%%%%%%%%%%%%%%%%%%%%%%%%%%%%%%%%%%%%%%%%%%%%%%%%%%%%%%%%%%

Complete pooling ignores group differences and uses one common parameter:

\[
\boxed{
\theta_1
=
\cdots
=
\theta_J
=
\theta.
}
\]

This can be too restrictive when genuine group heterogeneity exists.

%%%%%%%%%%%%%%%%%%%%%%%%%%%%%%%%%%%%%%%%%%%%%%%%%%%%%%%%%%%%
\subsection{No Pooling}
\label{subsec:no-pooling}
%%%%%%%%%%%%%%%%%%%%%%%%%%%%%%%%%%%%%%%%%%%%%%%%%%%%%%%%%%%%

No pooling estimates each group independently.

This avoids forcing group equality but can produce unstable estimates for
small groups.

Hierarchical partial pooling borrows information across groups while
allowing genuine differences.

%%%%%%%%%%%%%%%%%%%%%%%%%%%%%%%%%%%%%%%%%%%%%%%%%%%%%%%%%%%%
\subsection{Exchangeability}
\label{subsec:bayesian-exchangeability}
%%%%%%%%%%%%%%%%%%%%%%%%%%%%%%%%%%%%%%%%%%%%%%%%%%%%%%%%%%%%

A collection of random variables is exchangeable if its joint
distribution is invariant under permutations.

For finite \(n\),

\[
\boxed{
p
(
x_1,\ldots,x_n
)
=
p
(
x_{\pi(1)},\ldots,x_{\pi(n)}
)
}
\]

for every permutation \(\pi\).

Exchangeability often provides the conceptual basis for hierarchical
modeling of similar groups.

%%%%%%%%%%%%%%%%%%%%%%%%%%%%%%%%%%%%%%%%%%%%%%%%%%%%%%%%%%%%
\subsection{Conditional Independence in Hierarchical Models}
\label{subsec:hierarchical-conditional-independence}
%%%%%%%%%%%%%%%%%%%%%%%%%%%%%%%%%%%%%%%%%%%%%%%%%%%%%%%%%%%%

Observations may be dependent marginally but independent conditional on a
shared latent parameter.

For example,

\[
Y_{1j},
\ldots,
Y_{n_jj}
\]

may be conditionally independent given

\[
\theta_j.
\]

After integrating over uncertainty in \(\theta_j\), observations within
the same group become correlated.

%%%%%%%%%%%%%%%%%%%%%%%%%%%%%%%%%%%%%%%%%%%%%%%%%%%%%%%%%%%%
\subsection{Empirical Bayes}
\label{subsec:empirical-bayes}
%%%%%%%%%%%%%%%%%%%%%%%%%%%%%%%%%%%%%%%%%%%%%%%%%%%%%%%%%%%%

Empirical Bayes methods estimate prior hyperparameters from the observed
data.

For example, instead of assigning a prior to

\[
(\mu,\tau^2),
\]

one may estimate them from the collection of groups and then compute
group-level posteriors.

This can produce useful shrinkage estimates but is not fully Bayesian
because hyperparameter uncertainty may be only partially represented.

%%%%%%%%%%%%%%%%%%%%%%%%%%%%%%%%%%%%%%%%%%%%%%%%%%%%%%%%%%%%
\subsection{Bayesian Linear Regression}
\label{subsec:bayesian-linear-regression}
%%%%%%%%%%%%%%%%%%%%%%%%%%%%%%%%%%%%%%%%%%%%%%%%%%%%%%%%%%%%

Consider

\[
\boxed{
\mathbf{y}
=
X\boldsymbol{\beta}
+
\boldsymbol{\varepsilon},
}
\]

with

\[
\boldsymbol{\varepsilon}
\sim
N
(
\mathbf{0},
\sigma^2I
).
\]

Bayesian regression assigns prior distributions to

\[
\boldsymbol{\beta}
\]

and possibly

\[
\sigma^2.
\]

%%%%%%%%%%%%%%%%%%%%%%%%%%%%%%%%%%%%%%%%%%%%%%%%%%%%%%%%%%%%
\subsection{Normal Prior for Regression Coefficients}
\label{subsec:normal-prior-regression}
%%%%%%%%%%%%%%%%%%%%%%%%%%%%%%%%%%%%%%%%%%%%%%%%%%%%%%%%%%%%

Suppose

\[
\boxed{
\boldsymbol{\beta}
\sim
N
(
\mathbf{m}_0,
V_0
).
}
\]

With known \(\sigma^2\), the posterior is normal.

The posterior precision matrix is

\[
\boxed{
V_n^{-1}
=
V_0^{-1}
+
\frac1{\sigma^2}
X^TX.
}
\]

%%%%%%%%%%%%%%%%%%%%%%%%%%%%%%%%%%%%%%%%%%%%%%%%%%%%%%%%%%%%
\subsection{Bayesian Regression Posterior Mean}
\label{subsec:bayesian-regression-posterior-mean}
%%%%%%%%%%%%%%%%%%%%%%%%%%%%%%%%%%%%%%%%%%%%%%%%%%%%%%%%%%%%

The posterior covariance is

\[
\boxed{
V_n
=
\left[
V_0^{-1}
+
\frac1{\sigma^2}
X^TX
\right]^{-1}.
}
\]

The posterior mean is

\[
\boxed{
\mathbf{m}_n
=
V_n
\left[
V_0^{-1}\mathbf{m}_0
+
\frac1{\sigma^2}
X^T\mathbf{y}
\right].
}
\]

Thus the posterior combines prior information with least-squares
information.

%%%%%%%%%%%%%%%%%%%%%%%%%%%%%%%%%%%%%%%%%%%%%%%%%%%%%%%%%%%%
\subsection{Bayesian Regression and Regularization}
\label{subsec:bayesian-regularization}
%%%%%%%%%%%%%%%%%%%%%%%%%%%%%%%%%%%%%%%%%%%%%%%%%%%%%%%%%%%%

Suppose

\[
\boldsymbol{\beta}
\sim
N
(
\mathbf{0},
\tau^2I
).
\]

The posterior mode minimizes an objective of the form

\[
\boxed{
\|\mathbf{y}-X\boldsymbol{\beta}\|_2^2
+
\lambda
\|\boldsymbol{\beta}\|_2^2,
}
\]

for a suitable

\[
\lambda.
\]

Thus Gaussian priors connect Bayesian regression to ridge
regularization.

%%%%%%%%%%%%%%%%%%%%%%%%%%%%%%%%%%%%%%%%%%%%%%%%%%%%%%%%%%%%
\subsection{Laplace Priors and Lasso Connection}
\label{subsec:laplace-prior-lasso}
%%%%%%%%%%%%%%%%%%%%%%%%%%%%%%%%%%%%%%%%%%%%%%%%%%%%%%%%%%%%

A Laplace prior has density proportional to

\[
\boxed{
e^{-\lambda|\beta|}.
}
\]

Using independent Laplace priors on regression coefficients produces a
posterior mode related to the lasso objective

\[
\boxed{
\|\mathbf{y}-X\boldsymbol{\beta}\|_2^2
+
\lambda
\sum_j|\beta_j|.
}
\]

Thus regularization penalties can often be interpreted as negative
log-priors.

%%%%%%%%%%%%%%%%%%%%%%%%%%%%%%%%%%%%%%%%%%%%%%%%%%%%%%%%%%%%
\subsection{Bayesian Logistic Regression}
\label{subsec:bayesian-logistic-regression}
%%%%%%%%%%%%%%%%%%%%%%%%%%%%%%%%%%%%%%%%%%%%%%%%%%%%%%%%%%%%

For binary responses,

\[
Y_i
\sim
\operatorname{Bernoulli}(p_i),
\]

with

\[
\boxed{
\operatorname{logit}(p_i)
=
\mathbf{x}_i^T
\boldsymbol{\beta}.
}
\]

Bayesian logistic regression places a prior on

\[
\boldsymbol{\beta}.
\]

The posterior typically has no simple closed form.

Numerical methods are therefore required.

%%%%%%%%%%%%%%%%%%%%%%%%%%%%%%%%%%%%%%%%%%%%%%%%%%%%%%%%%%%%
\subsection{Why Computation Is Central}
\label{subsec:why-bayesian-computation}
%%%%%%%%%%%%%%%%%%%%%%%%%%%%%%%%%%%%%%%%%%%%%%%%%%%%%%%%%%%%

Many Bayesian posterior distributions involve integrals such as

\[
\boxed{
\int
g(\theta)
\pi(\theta\mid x)
\,d\theta
}
\]

that cannot be evaluated analytically.

Modern Bayesian statistics therefore relies heavily on:

\[
\boxed{
\text{Monte Carlo methods},
}
\]

\[
\boxed{
\text{numerical optimization},
}
\]

and

\[
\boxed{
\text{approximation algorithms}.
}
\]

%%%%%%%%%%%%%%%%%%%%%%%%%%%%%%%%%%%%%%%%%%%%%%%%%%%%%%%%%%%%
\subsection{Monte Carlo Integration}
\label{subsec:monte-carlo-integration-bayes}
%%%%%%%%%%%%%%%%%%%%%%%%%%%%%%%%%%%%%%%%%%%%%%%%%%%%%%%%%%%%

Suppose

\[
\theta^{(1)},
\ldots,
\theta^{(M)}
\]

are independent draws from the posterior.

Then

\[
\boxed{
\mathbb{E}
[
g(\theta)
\mid x
]
\approx
\frac1M
\sum_{m=1}^{M}
g
(
\theta^{(m)}
).
}
\]

This is ordinary Monte Carlo integration.

%%%%%%%%%%%%%%%%%%%%%%%%%%%%%%%%%%%%%%%%%%%%%%%%%%%%%%%%%%%%
\subsection{Worked Example: Monte Carlo Posterior Mean}
\label{subsec:worked-monte-carlo-posterior}
%%%%%%%%%%%%%%%%%%%%%%%%%%%%%%%%%%%%%%%%%%%%%%%%%%%%%%%%%%%%

\begin{workedexamplebox}{Posterior Simulation}

Suppose posterior draws are

\[
2.1,\ 2.5,\ 1.9,\ 2.7,\ 2.3.
\]

The Monte Carlo estimate of the posterior mean is

\[
\frac{
2.1+2.5+1.9+2.7+2.3
}{5}.
\]

Therefore,

\begin{answerbox}
\[
\boxed{
\widehat{\mathbb{E}}
[
\theta\mid x
]
=
2.3.
}
\]
\end{answerbox}

\end{workedexamplebox}

%%%%%%%%%%%%%%%%%%%%%%%%%%%%%%%%%%%%%%%%%%%%%%%%%%%%%%%%%%%%
\subsection{Markov Chain Monte Carlo}
\label{subsec:markov-chain-monte-carlo}
%%%%%%%%%%%%%%%%%%%%%%%%%%%%%%%%%%%%%%%%%%%%%%%%%%%%%%%%%%%%

Markov chain Monte Carlo, abbreviated MCMC, constructs a Markov chain

\[
\theta^{(1)},
\theta^{(2)},
\ldots
\]

whose stationary distribution is the posterior distribution.

After suitable convergence,

\[
\boxed{
\theta^{(m)}
\approx
\text{posterior draws}.
}
\]

These draws can be used to approximate posterior expectations and
probabilities.

%%%%%%%%%%%%%%%%%%%%%%%%%%%%%%%%%%%%%%%%%%%%%%%%%%%%%%%%%%%%
\subsection{Metropolis--Hastings Algorithm}
\label{subsec:metropolis-hastings}
%%%%%%%%%%%%%%%%%%%%%%%%%%%%%%%%%%%%%%%%%%%%%%%%%%%%%%%%%%%%

Suppose the current state is

\[
\theta.
\]

Propose a candidate

\[
\theta'
\sim
q(\theta'\mid\theta).
\]

Accept the proposal with probability

\[
\boxed{
\alpha
=
\min
\left\{
1,
\frac{
\pi(\theta'\mid x)
q(\theta\mid\theta')
}{
\pi(\theta\mid x)
q(\theta'\mid\theta)
}
\right\}.
}
\]

If rejected, retain the current state.

%%%%%%%%%%%%%%%%%%%%%%%%%%%%%%%%%%%%%%%%%%%%%%%%%%%%%%%%%%%%
\subsection{Why the Normalizing Constant Cancels}
\label{subsec:mcmc-normalizing-constant}
%%%%%%%%%%%%%%%%%%%%%%%%%%%%%%%%%%%%%%%%%%%%%%%%%%%%%%%%%%%%

Suppose

\[
\pi(\theta\mid x)
=
\frac{
h(\theta)
}{
C
}.
\]

In the Metropolis--Hastings ratio,

\[
\frac{
\pi(\theta'\mid x)
}{
\pi(\theta\mid x)
}
=
\frac{
h(\theta')/C
}{
h(\theta)/C
}
=
\frac{
h(\theta')
}{
h(\theta)
}.
\]

Thus the unknown normalizing constant cancels.

This is one major reason MCMC is powerful for complicated posteriors.

%%%%%%%%%%%%%%%%%%%%%%%%%%%%%%%%%%%%%%%%%%%%%%%%%%%%%%%%%%%%
\subsection{Random-Walk Metropolis}
\label{subsec:random-walk-metropolis}
%%%%%%%%%%%%%%%%%%%%%%%%%%%%%%%%%%%%%%%%%%%%%%%%%%%%%%%%%%%%

A common proposal is

\[
\boxed{
\theta'
=
\theta+\varepsilon,
}
\]

where

\[
\varepsilon
\]

is drawn from a symmetric distribution such as

\[
N(0,s^2).
\]

Then

\[
q(\theta'\mid\theta)
=
q(\theta\mid\theta'),
\]

so the proposal terms cancel.

The acceptance probability becomes

\[
\boxed{
\alpha
=
\min
\left\{
1,
\frac{
\pi(\theta'\mid x)
}{
\pi(\theta\mid x)
}
\right\}.
}
\]

%%%%%%%%%%%%%%%%%%%%%%%%%%%%%%%%%%%%%%%%%%%%%%%%%%%%%%%%%%%%
\subsection{Gibbs Sampling}
\label{subsec:gibbs-sampling}
%%%%%%%%%%%%%%%%%%%%%%%%%%%%%%%%%%%%%%%%%%%%%%%%%%%%%%%%%%%%

Suppose the posterior involves parameters

\[
\theta_1,\ldots,\theta_k.
\]

Gibbs sampling repeatedly draws from full conditional distributions:

\[
\boxed{
\theta_1
\sim
\pi
(
\theta_1
\mid
\theta_2,\ldots,\theta_k,x
),
}
\]

then

\[
\boxed{
\theta_2
\sim
\pi
(
\theta_2
\mid
\theta_1,\theta_3,\ldots,\theta_k,x
),
}
\]

and so forth.

When these conditionals are easy to sample, Gibbs sampling can be highly
effective.

%%%%%%%%%%%%%%%%%%%%%%%%%%%%%%%%%%%%%%%%%%%%%%%%%%%%%%%%%%%%
\subsection{Full Conditional Distribution}
\label{subsec:full-conditional-distribution}
%%%%%%%%%%%%%%%%%%%%%%%%%%%%%%%%%%%%%%%%%%%%%%%%%%%%%%%%%%%%

The full conditional for parameter \(\theta_j\) is

\[
\boxed{
\pi
(
\theta_j
\mid
\theta_{-j},x
),
}
\]

where

\[
\theta_{-j}
\]

denotes all parameters except \(\theta_j\).

In many hierarchical models, conjugacy makes full conditionals belong to
standard probability families.

%%%%%%%%%%%%%%%%%%%%%%%%%%%%%%%%%%%%%%%%%%%%%%%%%%%%%%%%%%%%
\subsection{Hamiltonian Monte Carlo Preview}
\label{subsec:hamiltonian-monte-carlo-preview}
%%%%%%%%%%%%%%%%%%%%%%%%%%%%%%%%%%%%%%%%%%%%%%%%%%%%%%%%%%%%

Hamiltonian Monte Carlo, abbreviated HMC, uses gradient information to
propose distant states with relatively high acceptance probability.

It introduces auxiliary momentum variables and simulates approximate
Hamiltonian dynamics.

HMC is particularly useful for high-dimensional continuous posterior
distributions.

%%%%%%%%%%%%%%%%%%%%%%%%%%%%%%%%%%%%%%%%%%%%%%%%%%%%%%%%%%%%
\subsection{No-U-Turn Sampler Preview}
\label{subsec:nuts-preview}
%%%%%%%%%%%%%%%%%%%%%%%%%%%%%%%%%%%%%%%%%%%%%%%%%%%%%%%%%%%%

The No-U-Turn Sampler, abbreviated NUTS, is an adaptive extension of HMC.

It automatically chooses trajectory lengths rather than requiring a fixed
number of leapfrog steps.

This reduces some tuning burden while preserving HMC's efficiency.

%%%%%%%%%%%%%%%%%%%%%%%%%%%%%%%%%%%%%%%%%%%%%%%%%%%%%%%%%%%%
\subsection{MCMC Autocorrelation}
\label{subsec:mcmc-autocorrelation}
%%%%%%%%%%%%%%%%%%%%%%%%%%%%%%%%%%%%%%%%%%%%%%%%%%%%%%%%%%%%

MCMC draws are generally not independent.

Successive states may be highly correlated.

Therefore

\[
M
\]

MCMC iterations do not necessarily contain the same information as

\[
M
\]

independent posterior draws.

%%%%%%%%%%%%%%%%%%%%%%%%%%%%%%%%%%%%%%%%%%%%%%%%%%%%%%%%%%%%
\subsection{Effective Sample Size}
\label{subsec:mcmc-effective-sample-size}
%%%%%%%%%%%%%%%%%%%%%%%%%%%%%%%%%%%%%%%%%%%%%%%%%%%%%%%%%%%%

The effective sample size, abbreviated ESS, measures the approximate
number of independent draws represented by a correlated Markov chain.

Conceptually,

\[
\boxed{
ESS
<
M
}
\]

when positive autocorrelation is present.

Higher ESS means more precise Monte Carlo estimates.

%%%%%%%%%%%%%%%%%%%%%%%%%%%%%%%%%%%%%%%%%%%%%%%%%%%%%%%%%%%%
\subsection{Monte Carlo Standard Error}
\label{subsec:monte-carlo-standard-error}
%%%%%%%%%%%%%%%%%%%%%%%%%%%%%%%%%%%%%%%%%%%%%%%%%%%%%%%%%%%%

Monte Carlo computation introduces simulation error in addition to
posterior uncertainty.

For an estimated posterior quantity,

\[
\boxed{
MCSE
}
\]

measures numerical uncertainty due to finite simulation.

A posterior analysis should use enough effective draws that MCSE is small
relative to the inferential uncertainty being reported.

%%%%%%%%%%%%%%%%%%%%%%%%%%%%%%%%%%%%%%%%%%%%%%%%%%%%%%%%%%%%
\subsection{Trace Plots}
\label{subsec:mcmc-trace-plots}
%%%%%%%%%%%%%%%%%%%%%%%%%%%%%%%%%%%%%%%%%%%%%%%%%%%%%%%%%%%%

A trace plot graphs sampled parameter values against iteration number.

A well-mixing chain should explore the posterior region repeatedly
without long trends or persistent trapping.

Trace plots help reveal:

\[
\boxed{
\text{poor mixing},
}
\]

\[
\boxed{
\text{nonstationarity},
}
\]

or

\[
\boxed{
\text{multimodality}.
}
\]

%%%%%%%%%%%%%%%%%%%%%%%%%%%%%%%%%%%%%%%%%%%%%%%%%%%%%%%%%%%%
\subsection{Multiple Chains}
\label{subsec:mcmc-multiple-chains}
%%%%%%%%%%%%%%%%%%%%%%%%%%%%%%%%%%%%%%%%%%%%%%%%%%%%%%%%%%%%

Running multiple chains from dispersed initial values helps assess
whether different simulations reach similar posterior regions.

Chains that remain separated may indicate:

\[
\boxed{
\text{lack of convergence},
}
\]

or

\[
\boxed{
\text{multimodal posterior structure}.
}
\]

%%%%%%%%%%%%%%%%%%%%%%%%%%%%%%%%%%%%%%%%%%%%%%%%%%%%%%%%%%%%
\subsection{Potential Scale Reduction}
\label{subsec:rhat-preview}
%%%%%%%%%%%%%%%%%%%%%%%%%%%%%%%%%%%%%%%%%%%%%%%%%%%%%%%%%%%%

A common convergence diagnostic compares within-chain and between-chain
variation.

It is often reported as

\[
\boxed{
\widehat{R}.
}
\]

Values sufficiently close to \(1\) indicate that between-chain and
within-chain behavior are similar.

No single threshold or diagnostic guarantees convergence.

%%%%%%%%%%%%%%%%%%%%%%%%%%%%%%%%%%%%%%%%%%%%%%%%%%%%%%%%%%%%
\subsection{Burn-In and Warmup}
\label{subsec:mcmc-warmup}
%%%%%%%%%%%%%%%%%%%%%%%%%%%%%%%%%%%%%%%%%%%%%%%%%%%%%%%%%%%%

Early MCMC iterations may reflect initialization or algorithm adaptation
rather than the target stationary distribution.

These iterations may be treated as warmup.

Modern samplers often use warmup to tune algorithm parameters.

Discarding early samples does not repair fundamentally poor convergence.

%%%%%%%%%%%%%%%%%%%%%%%%%%%%%%%%%%%%%%%%%%%%%%%%%%%%%%%%%%%%
\subsection{Thinning}
\label{subsec:mcmc-thinning}
%%%%%%%%%%%%%%%%%%%%%%%%%%%%%%%%%%%%%%%%%%%%%%%%%%%%%%%%%%%%

Thinning retains only every \(k\)-th MCMC draw.

Although this reduces storage and autocorrelation among retained draws,
it also discards information.

For most modern analyses, improving the sampler or retaining all valid
post-warmup draws is preferable unless storage or specific computational
constraints justify thinning.

%%%%%%%%%%%%%%%%%%%%%%%%%%%%%%%%%%%%%%%%%%%%%%%%%%%%%%%%%%%%
\subsection{Importance Sampling Preview}
\label{subsec:importance-sampling-bayes}
%%%%%%%%%%%%%%%%%%%%%%%%%%%%%%%%%%%%%%%%%%%%%%%%%%%%%%%%%%%%

Suppose direct simulation from target density

\[
\pi(\theta)
\]

is difficult.

Choose proposal density

\[
q(\theta)
\]

and define weights

\[
\boxed{
w(\theta)
=
\frac{
\pi(\theta)
}{
q(\theta)
}.
}
\]

Expectations under \(\pi\) can then be approximated using weighted draws
from \(q\).

Poor overlap between \(q\) and \(\pi\) can cause unstable weights.

%%%%%%%%%%%%%%%%%%%%%%%%%%%%%%%%%%%%%%%%%%%%%%%%%%%%%%%%%%%%
\subsection{Laplace Approximation Preview}
\label{subsec:laplace-approximation}
%%%%%%%%%%%%%%%%%%%%%%%%%%%%%%%%%%%%%%%%%%%%%%%%%%%%%%%%%%%%

Suppose the posterior is concentrated near mode

\[
\widehat{\theta}.
\]

A second-order Taylor expansion of the log posterior yields an approximate
normal distribution:

\[
\boxed{
\theta\mid x
\approx
N
\left(
\widehat{\theta},
H^{-1}
\right),
}
\]

where \(H\) is the negative Hessian of the log posterior at the mode.

This is the Laplace approximation.

%%%%%%%%%%%%%%%%%%%%%%%%%%%%%%%%%%%%%%%%%%%%%%%%%%%%%%%%%%%%
\subsection{Variational Inference Preview}
\label{subsec:variational-inference-preview}
%%%%%%%%%%%%%%%%%%%%%%%%%%%%%%%%%%%%%%%%%%%%%%%%%%%%%%%%%%%%

Variational inference approximates a difficult posterior

\[
\pi(\theta\mid x)
\]

with a simpler distribution

\[
q_\phi(\theta)
\]

chosen from a tractable family.

The parameters \(\phi\) are selected to make \(q_\phi\) close to the
posterior according to an optimization criterion.

Thus Bayesian computation can be converted into an optimization problem.

%%%%%%%%%%%%%%%%%%%%%%%%%%%%%%%%%%%%%%%%%%%%%%%%%%%%%%%%%%%%
\subsection{Kullback--Leibler Divergence Preview}
\label{subsec:kl-divergence-bayes}
%%%%%%%%%%%%%%%%%%%%%%%%%%%%%%%%%%%%%%%%%%%%%%%%%%%%%%%%%%%%

A common discrepancy is the Kullback--Leibler divergence:

\[
\boxed{
D_{KL}(q\|p)
=
\int
q(\theta)
\log
\frac{
q(\theta)
}{
p(\theta)
}
\,d\theta.
}
\]

Variational methods often minimize

\[
D_{KL}
(
q_\phi
\|
\pi(\cdot\mid x)
).
\]

KL divergence is not symmetric and is not a metric.

%%%%%%%%%%%%%%%%%%%%%%%%%%%%%%%%%%%%%%%%%%%%%%%%%%%%%%%%%%%%
\subsection{Evidence Lower Bound Preview}
\label{subsec:elbo-preview}
%%%%%%%%%%%%%%%%%%%%%%%%%%%%%%%%%%%%%%%%%%%%%%%%%%%%%%%%%%%%

Variational inference commonly maximizes the evidence lower bound,
abbreviated ELBO.

A standard form is

\[
\boxed{
ELBO
=
\mathbb{E}_q
[
\log p(x,\theta)
]
-
\mathbb{E}_q
[
\log q(\theta)
].
}
\]

Maximizing the ELBO is equivalent to minimizing

\[
D_{KL}
(
q
\|
\pi(\theta\mid x)
)
\]

up to the constant

\[
\log p(x).
\]

%%%%%%%%%%%%%%%%%%%%%%%%%%%%%%%%%%%%%%%%%%%%%%%%%%%%%%%%%%%%
\subsection{Sensitivity Analysis}
\label{subsec:bayesian-sensitivity-analysis}
%%%%%%%%%%%%%%%%%%%%%%%%%%%%%%%%%%%%%%%%%%%%%%%%%%%%%%%%%%%%

Bayesian conclusions may depend on prior choices.

A sensitivity analysis refits or re-evaluates the model under several
scientifically plausible priors.

One may compare:

\[
\boxed{
\text{posterior means},
}
\]

\[
\boxed{
\text{credible intervals},
}
\]

\[
\boxed{
\text{predictive distributions},
}
\]

or

\[
\boxed{
\text{decision conclusions}.
}
\]

%%%%%%%%%%%%%%%%%%%%%%%%%%%%%%%%%%%%%%%%%%%%%%%%%%%%%%%%%%%%
\subsection{Prior Dominance}
\label{subsec:prior-dominance}
%%%%%%%%%%%%%%%%%%%%%%%%%%%%%%%%%%%%%%%%%%%%%%%%%%%%%%%%%%%%

If data are weak and the prior is highly concentrated, the posterior may
remain close to the prior.

This is not automatically a flaw if the prior information is genuine.

However, the analysis should clearly communicate the relative
contribution of prior and data information.

%%%%%%%%%%%%%%%%%%%%%%%%%%%%%%%%%%%%%%%%%%%%%%%%%%%%%%%%%%%%
\subsection{Likelihood Dominance}
\label{subsec:likelihood-dominance}
%%%%%%%%%%%%%%%%%%%%%%%%%%%%%%%%%%%%%%%%%%%%%%%%%%%%%%%%%%%%

With highly informative data, the likelihood may dominate the prior over
the posterior-relevant region.

Then several reasonable priors may lead to similar posterior
distributions.

This is one form of Bayesian robustness to prior specification in large
samples.

%%%%%%%%%%%%%%%%%%%%%%%%%%%%%%%%%%%%%%%%%%%%%%%%%%%%%%%%%%%%
\subsection{Bayesian Hypothesis Testing}
\label{subsec:bayesian-hypothesis-testing}
%%%%%%%%%%%%%%%%%%%%%%%%%%%%%%%%%%%%%%%%%%%%%%%%%%%%%%%%%%%%

Bayesian testing may compare hypotheses through:

\[
\boxed{
\text{posterior probabilities},
}
\]

\[
\boxed{
\text{posterior odds},
}
\]

or

\[
\boxed{
\text{Bayes factors}.
}
\]

This differs fundamentally from frequentist \(p\)-values.

A posterior probability such as

\[
P(H_0\mid x)
\]

is not a \(p\)-value.

%%%%%%%%%%%%%%%%%%%%%%%%%%%%%%%%%%%%%%%%%%%%%%%%%%%%%%%%%%%%
\subsection{Point Null Hypotheses}
\label{subsec:point-null-bayes}
%%%%%%%%%%%%%%%%%%%%%%%%%%%%%%%%%%%%%%%%%%%%%%%%%%%%%%%%%%%%

If one wishes to assign positive posterior probability to an exact null

\[
H_0:\theta=\theta_0,
\]

the prior model must assign positive prior probability to that point.

For example, one may use a mixture prior

\[
\boxed{
\pi(\theta)
=
w\Delta_{\theta_0}
+
(1-w)\pi_1(\theta),
}
\]

where

\[
\Delta_{\theta_0}
\]

is a point mass at \(\theta_0\).

%%%%%%%%%%%%%%%%%%%%%%%%%%%%%%%%%%%%%%%%%%%%%%%%%%%%%%%%%%%%
\subsection{Continuous Priors and Exact Point Probabilities}
\label{subsec:continuous-prior-point-null}
%%%%%%%%%%%%%%%%%%%%%%%%%%%%%%%%%%%%%%%%%%%%%%%%%%%%%%%%%%%%

If the prior for a continuous parameter has a continuous density only,
then

\[
\boxed{
P(\theta=\theta_0)=0.
}
\]

The posterior also assigns zero probability to that exact point.

Thus testing an exact point null requires a model structure different
from merely examining a continuous posterior density.

%%%%%%%%%%%%%%%%%%%%%%%%%%%%%%%%%%%%%%%%%%%%%%%%%%%%%%%%%%%%
\subsection{Posterior Probability of an Interval}
\label{subsec:posterior-probability-interval}
%%%%%%%%%%%%%%%%%%%%%%%%%%%%%%%%%%%%%%%%%%%%%%%%%%%%%%%%%%%%

Instead of testing exact equality, one may ask whether \(\theta\) lies in
a scientifically meaningful region:

\[
\boxed{
P
(
a<\theta<b
\mid x
).
}
\]

This can directly quantify evidence that an effect is practically small,
positive, negative, or clinically meaningful.

%%%%%%%%%%%%%%%%%%%%%%%%%%%%%%%%%%%%%%%%%%%%%%%%%%%%%%%%%%%%
\subsection{Region of Practical Equivalence Preview}
\label{subsec:rope-preview}
%%%%%%%%%%%%%%%%%%%%%%%%%%%%%%%%%%%%%%%%%%%%%%%%%%%%%%%%%%%%

A region of practical equivalence, sometimes abbreviated ROPE, specifies
parameter values considered practically negligible.

For example,

\[
\boxed{
|\theta|<\delta.
}
\]

The posterior probability

\[
\boxed{
P
(
|\theta|<\delta
\mid x
)
}
\]

then measures support for a practically small effect.

The choice of \(\delta\) must be scientifically justified.

%%%%%%%%%%%%%%%%%%%%%%%%%%%%%%%%%%%%%%%%%%%%%%%%%%%%%%%%%%%%
\subsection{Posterior Predictive Decision Making}
\label{subsec:posterior-predictive-decision}
%%%%%%%%%%%%%%%%%%%%%%%%%%%%%%%%%%%%%%%%%%%%%%%%%%%%%%%%%%%%

Many practical decisions depend more directly on future outcomes than on
parameters.

One can therefore evaluate expected utility or loss under the posterior
predictive distribution:

\[
\boxed{
p
(
y_{\text{new}}
\mid x
).
}
\]

This directly connects Bayesian inference with forecasting and decision
analysis.

%%%%%%%%%%%%%%%%%%%%%%%%%%%%%%%%%%%%%%%%%%%%%%%%%%%%%%%%%%%%
\subsection{Bayesian Model Averaging Preview}
\label{subsec:bayesian-model-averaging}
%%%%%%%%%%%%%%%%%%%%%%%%%%%%%%%%%%%%%%%%%%%%%%%%%%%%%%%%%%%%

Suppose several models

\[
M_1,\ldots,M_K
\]

have posterior probabilities.

Rather than choosing one model, Bayesian model averaging predicts by

\[
\boxed{
p
(
y_{\text{new}}
\mid x
)
=
\sum_{k=1}^{K}
p
(
y_{\text{new}}
\mid x,M_k
)
P
(
M_k\mid x
).
}
\]

Thus model uncertainty is propagated into prediction.

%%%%%%%%%%%%%%%%%%%%%%%%%%%%%%%%%%%%%%%%%%%%%%%%%%%%%%%%%%%%
\subsection{Bayesian Workflow}
\label{subsec:bayesian-workflow}
%%%%%%%%%%%%%%%%%%%%%%%%%%%%%%%%%%%%%%%%%%%%%%%%%%%%%%%%%%%%

A practical Bayesian workflow is:

\begin{enumerate}
    \item define the scientific estimand;
    \item specify the likelihood model;
    \item identify parameter constraints;
    \item choose prior distributions;
    \item perform prior predictive checks;
    \item derive or compute the posterior;
    \item check numerical convergence when simulation is used;
    \item summarize posterior quantities;
    \item construct credible intervals or posterior probabilities;
    \item compute posterior predictive distributions;
    \item perform posterior predictive model checks;
    \item evaluate prior sensitivity;
    \item compare models when scientifically necessary;
    \item make decisions using explicit losses or utilities when
          appropriate;
    \item report model assumptions and computational uncertainty.
\end{enumerate}

%%%%%%%%%%%%%%%%%%%%%%%%%%%%%%%%%%%%%%%%%%%%%%%%%%%%%%%%%%%%
\subsection{Choosing Priors}
\label{subsec:choosing-priors}
%%%%%%%%%%%%%%%%%%%%%%%%%%%%%%%%%%%%%%%%%%%%%%%%%%%%%%%%%%%%

A useful prior should reflect:

\[
\boxed{
\text{parameter support},
}
\]

\[
\boxed{
\text{scientifically plausible magnitudes},
}
\]

and

\[
\boxed{
\text{appropriate prior uncertainty}.
}
\]

Priors should not be chosen solely because they produce convenient
posterior conclusions.

%%%%%%%%%%%%%%%%%%%%%%%%%%%%%%%%%%%%%%%%%%%%%%%%%%%%%%%%%%%%
\subsection{Common Errors}
\label{subsec:bayesian-common-errors}
%%%%%%%%%%%%%%%%%%%%%%%%%%%%%%%%%%%%%%%%%%%%%%%%%%%%%%%%%%%%

\subsubsection{Treating the Likelihood as a Probability Distribution Over the Parameter}

The likelihood is proportional to the sampling probability of the
observed data as a function of the parameter.

It does not integrate to one over \(\theta\) in general.

%%%%%%%%%%%%%%%%%%%%%%%%%%%%%%%%%%%%%%%%%%%%%%%%%%%%%%%%%%%%
\subsubsection{Forgetting the Normalizing Constant}

When identifying posterior families, proportionality is sufficient.

But final posterior probabilities and densities require a properly
normalized distribution.

%%%%%%%%%%%%%%%%%%%%%%%%%%%%%%%%%%%%%%%%%%%%%%%%%%%%%%%%%%%%
\subsubsection{Confusing Prior and Posterior Information}

The prior describes uncertainty before the current data.

The posterior combines that prior with the current likelihood.

%%%%%%%%%%%%%%%%%%%%%%%%%%%%%%%%%%%%%%%%%%%%%%%%%%%%%%%%%%%%
\subsubsection{Interpreting a Credible Interval as a Frequentist Confidence Interval}

The two intervals may look numerically similar but have different
probability interpretations.

%%%%%%%%%%%%%%%%%%%%%%%%%%%%%%%%%%%%%%%%%%%%%%%%%%%%%%%%%%%%
\subsubsection{Interpreting a Confidence Interval as Having Posterior Probability}

A frequentist \(95\%\) confidence interval does not generally imply

\[
P(\theta\text{ lies in the observed interval})=0.95
\]

under the frequentist framework.

%%%%%%%%%%%%%%%%%%%%%%%%%%%%%%%%%%%%%%%%%%%%%%%%%%%%%%%%%%%%
\subsubsection{Calling a Prior Noninformative Without Checking Parameterization}

A prior that is flat in \(\theta\) need not be flat in

\[
g(\theta).
\]

Thus ``noninformative'' depends on the parameter representation.

%%%%%%%%%%%%%%%%%%%%%%%%%%%%%%%%%%%%%%%%%%%%%%%%%%%%%%%%%%%%
\subsubsection{Using Improper Priors Without Checking Posterior Propriety}

An improper posterior is not a valid probability distribution.

Posterior normalization must exist.

%%%%%%%%%%%%%%%%%%%%%%%%%%%%%%%%%%%%%%%%%%%%%%%%%%%%%%%%%%%%
\subsubsection{Using Improper Priors Naively for Bayes Factors}

Bayes factors depend on prior normalizing constants.

Improper priors can make ordinary marginal likelihoods undefined up to
arbitrary multiplicative constants.

%%%%%%%%%%%%%%%%%%%%%%%%%%%%%%%%%%%%%%%%%%%%%%%%%%%%%%%%%%%%
\subsubsection{Ignoring Prior Predictive Implications}

A mathematically broad prior may imply impossible or absurd datasets.

Prior predictive checks should be performed.

%%%%%%%%%%%%%%%%%%%%%%%%%%%%%%%%%%%%%%%%%%%%%%%%%%%%%%%%%%%%
\subsubsection{Assuming More Posterior Concentration Means the Model Is Correct}

A misspecified model can yield a very concentrated posterior around the
wrong parameter value.

Model adequacy must be checked separately.

%%%%%%%%%%%%%%%%%%%%%%%%%%%%%%%%%%%%%%%%%%%%%%%%%%%%%%%%%%%%
\subsubsection{Using MAP as Though It Were the Entire Posterior}

The MAP is only one point summary.

It does not communicate posterior uncertainty or skewness.

%%%%%%%%%%%%%%%%%%%%%%%%%%%%%%%%%%%%%%%%%%%%%%%%%%%%%%%%%%%%
\subsubsection{Reporting Only Posterior Means}

Posterior distributions may be skewed, multimodal, or heavy-tailed.

Credible intervals and graphical summaries may be essential.

%%%%%%%%%%%%%%%%%%%%%%%%%%%%%%%%%%%%%%%%%%%%%%%%%%%%%%%%%%%%
\subsubsection{Ignoring MCMC Dependence}

MCMC output is correlated.

The number of iterations is not the same as the effective number of
independent posterior samples.

%%%%%%%%%%%%%%%%%%%%%%%%%%%%%%%%%%%%%%%%%%%%%%%%%%%%%%%%%%%%
\subsubsection{Assuming a Good Trace Plot Proves Convergence}

Trace plots are useful but not definitive.

Multiple chains, effective sample sizes, convergence diagnostics, and
problem-specific checks should be considered together.

%%%%%%%%%%%%%%%%%%%%%%%%%%%%%%%%%%%%%%%%%%%%%%%%%%%%%%%%%%%%
\subsubsection{Thinning Automatically}

Thinning discards valid posterior draws and often reduces computational
efficiency.

It should be used for a specific reason rather than as a default.

%%%%%%%%%%%%%%%%%%%%%%%%%%%%%%%%%%%%%%%%%%%%%%%%%%%%%%%%%%%%
\subsubsection{Ignoring Posterior Predictive Performance}

A model may fit parameters well but produce poor predictions.

Posterior predictive checks evaluate the data-generating implications of
the fitted model.

%%%%%%%%%%%%%%%%%%%%%%%%%%%%%%%%%%%%%%%%%%%%%%%%%%%%%%%%%%%%
\subsubsection{Treating Bayes Factors as Prior-Free Evidence Measures}

Bayes factors depend on the prior distributions within the compared
models.

Prior sensitivity can be substantial.

%%%%%%%%%%%%%%%%%%%%%%%%%%%%%%%%%%%%%%%%%%%%%%%%%%%%%%%%%%%%
\subsubsection{Treating Posterior Probability as Scientific Importance}

A parameter may have high posterior probability of being positive while
the effect magnitude is practically negligible.

Magnitude and decision relevance must also be considered.

%%%%%%%%%%%%%%%%%%%%%%%%%%%%%%%%%%%%%%%%%%%%%%%%%%%%%%%%%%%%
\subsection{Applications}
\label{subsec:bayesian-applications}
%%%%%%%%%%%%%%%%%%%%%%%%%%%%%%%%%%%%%%%%%%%%%%%%%%%%%%%%%%%%

\begin{applicationbox}

\textbf{Scientific Research.}
Bayesian inference combines previous evidence with new observations and
produces probability distributions for parameters and predictions.

\medskip

\textbf{Medicine.}
Bayesian methods are used for clinical trials, dose finding, hierarchical
meta-analysis, diagnostic models, and adaptive decision making.

\medskip

\textbf{Engineering.}
Bayesian models update reliability, failure-rate, and system-state
estimates as new data become available.

\medskip

\textbf{Environmental Science.}
Hierarchical Bayesian models combine information across locations,
times, sensors, and ecological systems.

\medskip

\textbf{Education Research.}
Partial pooling can estimate school, class, or program effects while
accounting for differing sample sizes.

\medskip

\textbf{Economics and Finance.}
Bayesian models incorporate parameter uncertainty into forecasting,
portfolio decisions, and dynamic models.

\medskip

\textbf{Machine Learning.}
Bayesian methods provide probabilistic prediction, regularization,
uncertainty quantification, latent-variable models, and model averaging.

\medskip

\textbf{Decision Analysis.}
Posterior distributions can be combined with explicit costs and benefits
to select actions minimizing expected loss.

\end{applicationbox}

%%%%%%%%%%%%%%%%%%%%%%%%%%%%%%%%%%%%%%%%%%%%%%%%%%%%%%%%%%%%
\subsection{Exercises}
\label{subsec:bayesian-exercises}
%%%%%%%%%%%%%%%%%%%%%%%%%%%%%%%%%%%%%%%%%%%%%%%%%%%%%%%%%%%%

\begin{exercise}[1]
State Bayes' theorem for a continuous parameter.
\end{exercise}

\begin{exercise}[1]
Define a prior distribution.
\end{exercise}

\begin{exercise}[1]
Define a posterior distribution.
\end{exercise}

\begin{exercise}[1]
Define a marginal likelihood.
\end{exercise}

\begin{exercise}[1]
Define a conjugate prior.
\end{exercise}

\begin{exercise}[1]
Define a credible interval.
\end{exercise}

\begin{exercise}[1]
Define a posterior predictive distribution.
\end{exercise}

\begin{exercise}[1]
Define a Bayes factor.
\end{exercise}

\begin{exercise}[1]
Define a hyperparameter.
\end{exercise}

\begin{exercise}[1]
Define partial pooling.
\end{exercise}

\begin{exercise}[1]
Define an improper prior.
\end{exercise}

\begin{exercise}[1]
Define posterior propriety.
\end{exercise}

\begin{exercise}[2]
Suppose

\[
p
\sim
\operatorname{Beta}(3,4)
\]

and \(10\) Bernoulli trials contain \(7\) successes. Find the posterior
distribution.
\end{exercise}

\begin{exercise}[2]
For the previous posterior, compute the posterior mean.
\end{exercise}

\begin{exercise}[2]
Suppose

\[
p
\sim
\operatorname{Beta}(1,1)
\]

and \(n=20\) observations contain \(15\) successes. Find the posterior
distribution.
\end{exercise}

\begin{exercise}[2]
Find the posterior predictive probability of success for the previous
problem.
\end{exercise}

\begin{exercise}[2]
Suppose

\[
\lambda
\sim
\operatorname{Gamma}(2,3)
\]

in shape--rate form.

Observe \(n=4\) Poisson counts with total \(S=10\). Find the posterior
distribution.
\end{exercise}

\begin{exercise}[2]
Compute the posterior mean for the previous Gamma--Poisson model.
\end{exercise}

\begin{exercise}[2]
Suppose

\[
\mu
\sim
N(0,4),
\]

\[
\sigma^2=9,
\qquad
n=9,
\qquad
\overline{x}=3.
\]

Find the posterior precision.
\end{exercise}

\begin{exercise}[2]
Compute the posterior mean and variance for the previous problem.
\end{exercise}

\begin{exercise}[2]
If

\[
P(H_1)=0.3,
\qquad
P(H_0)=0.7,
\]

and

\[
BF_{10}=5,
\]

compute the posterior odds for \(H_1\) against \(H_0\).
\end{exercise}

\begin{exercise}[2]
Convert the posterior odds from the previous problem into a posterior
probability for \(H_1\).
\end{exercise}

\begin{exercise}[2]
Explain why the posterior mean is the Bayes estimator under squared-error
loss.
\end{exercise}

\begin{exercise}[2]
Explain why the posterior median is optimal under absolute-error loss.
\end{exercise}

\begin{exercise}[3]
Derive the Beta--Bernoulli posterior.
\end{exercise}

\begin{exercise}[3]
Show that the Beta posterior mean is a weighted average of the prior mean
and sample proportion.
\end{exercise}

\begin{exercise}[3]
Derive the Beta--Binomial posterior predictive distribution.
\end{exercise}

\begin{exercise}[3]
Derive the Gamma--Poisson posterior.
\end{exercise}

\begin{exercise}[3]
Derive the posterior predictive distribution for a future Poisson count
under a Gamma prior.
\end{exercise}

\begin{exercise}[3]
Derive the Normal--Normal posterior mean and variance.
\end{exercise}

\begin{exercise}[3]
Show that the posterior precision equals prior precision plus data
precision in the Normal--Normal model.
\end{exercise}

\begin{exercise}[3]
Explain why posterior means shrink toward prior means.
\end{exercise}

\begin{exercise}[3]
Compare equal-tailed and HPD credible intervals.
\end{exercise}

\begin{exercise}[3]
Explain why a \(95\%\) credible interval and a \(95\%\) confidence
interval have different interpretations.
\end{exercise}

\begin{exercise}[3]
Derive the posterior predictive distribution as an integral over the
posterior.
\end{exercise}

\begin{exercise}[3]
Explain the role of prior predictive checks.
\end{exercise}

\begin{exercise}[3]
Explain the role of posterior predictive checks.
\end{exercise}

\begin{exercise}[3]
Explain why improper priors can sometimes yield proper posteriors.
\end{exercise}

\begin{exercise}[3]
Explain why improper priors create problems for ordinary Bayes factors.
\end{exercise}

\begin{exercise}[3]
Derive posterior odds as prior odds multiplied by a Bayes factor.
\end{exercise}

\begin{exercise}[3]
Explain why marginal likelihood automatically averages over parameter
uncertainty.
\end{exercise}

\begin{exercise}[3]
Explain the difference between complete pooling, no pooling, and partial
pooling.
\end{exercise}

\begin{exercise}[3]
Explain how hierarchical models induce shrinkage.
\end{exercise}

\begin{exercise}[3]
Explain the role of exchangeability in hierarchical modeling.
\end{exercise}

\begin{exercise}[3]
Derive the Gaussian-prior Bayesian linear-regression posterior with known
variance.
\end{exercise}

\begin{exercise}[3]
Show the connection between Gaussian priors and ridge regression.
\end{exercise}

\begin{exercise}[3]
Show the connection between Laplace priors and lasso MAP estimation.
\end{exercise}

\begin{exercise}[3]
Explain why Bayesian logistic regression usually requires numerical
methods.
\end{exercise}

\begin{exercise}[3]
Explain why the normalizing constant cancels in Metropolis--Hastings.
\end{exercise}

\begin{exercise}[3]
Write the Metropolis--Hastings acceptance probability for a symmetric
proposal.
\end{exercise}

\begin{exercise}[3]
Explain the difference between independent Monte Carlo samples and MCMC
samples.
\end{exercise}

\begin{exercise}[3]
Explain why effective sample size is smaller than the raw iteration count
when autocorrelation is positive.
\end{exercise}

\begin{exercise}[3]
Explain why thinning is not generally required.
\end{exercise}

\begin{exercise}[3]
Explain how prior sensitivity analysis should be performed.
\end{exercise}

\begin{exercise}[4]
Derive the general conjugate prior for a regular exponential family.
\end{exercise}

\begin{exercise}[4]
Study Jeffreys priors derived from Fisher information.
\end{exercise}

\begin{exercise}[4]
Investigate reparameterization invariance of Jeffreys priors.
\end{exercise}

\begin{exercise}[4]
Study reference priors and objective Bayesian methods.
\end{exercise}

\begin{exercise}[4]
Derive the Normal--Inverse-Gamma posterior for a normal model with unknown
mean and variance.
\end{exercise}

\begin{exercise}[4]
Derive the multivariate Normal--Inverse-Wishart conjugate posterior.
\end{exercise}

\begin{exercise}[4]
Study posterior consistency.
\end{exercise}

\begin{exercise}[4]
Study the Bernstein--von Mises theorem.
\end{exercise}

\begin{exercise}[4]
Investigate Bayesian credible sets and frequentist asymptotic coverage.
\end{exercise}

\begin{exercise}[4]
Derive Bayes factors for simple versus simple hypotheses.
\end{exercise}

\begin{exercise}[4]
Study Bayes factors for normal mean testing.
\end{exercise}

\begin{exercise}[4]
Investigate the Jeffreys--Lindley phenomenon.
\end{exercise}

\begin{exercise}[4]
Study Bayesian model averaging.
\end{exercise}

\begin{exercise}[4]
Derive the posterior distribution in a normal hierarchical model.
\end{exercise}

\begin{exercise}[4]
Investigate empirical Bayes estimation.
\end{exercise}

\begin{exercise}[4]
Study de Finetti's theorem and infinite exchangeability.
\end{exercise}

\begin{exercise}[4]
Derive the Metropolis--Hastings detailed balance condition.
\end{exercise}

\begin{exercise}[4]
Study ergodicity conditions for MCMC.
\end{exercise}

\begin{exercise}[4]
Derive Gibbs sampling for a conjugate hierarchical model.
\end{exercise}

\begin{exercise}[4]
Study Hamiltonian Monte Carlo.
\end{exercise}

\begin{exercise}[4]
Investigate convergence diagnostics and effective sample size.
\end{exercise}

\begin{exercise}[4]
Study importance sampling variance and weight degeneracy.
\end{exercise}

\begin{exercise}[4]
Derive the Laplace approximation.
\end{exercise}

\begin{exercise}[4]
Derive the evidence lower bound for variational inference.
\end{exercise}

\begin{exercise}[4]
Study posterior predictive checks and Bayesian \(p\)-value concepts.
\end{exercise}

%%%%%%%%%%%%%%%%%%%%%%%%%%%%%%%%%%%%%%%%%%%%%%%%%%%%%%%%%%%%
\subsection{Conceptual Summary}
\label{subsec:bayesian-summary}
%%%%%%%%%%%%%%%%%%%%%%%%%%%%%%%%%%%%%%%%%%%%%%%%%%%%%%%%%%%%

Bayesian inference is based on

\[
\boxed{
\pi(\theta\mid x)
\propto
L(\theta;x)\pi(\theta).
}
\]

The prior is

\[
\boxed{
\pi(\theta),
}
\]

the likelihood is

\[
\boxed{
L(\theta;x),
}
\]

and the posterior is

\[
\boxed{
\pi(\theta\mid x).
}
\]

The marginal likelihood is

\[
\boxed{
m(x)
=
\int
L(\theta;x)
\pi(\theta)
\,d\theta.
}
\]

For Beta--Bernoulli updating,

\[
\boxed{
p\mid x
\sim
\operatorname{Beta}
(
\alpha+S,
\beta+n-S
).
}
\]

For Gamma--Poisson updating,

\[
\boxed{
\lambda\mid x
\sim
\operatorname{Gamma}
(
\alpha+S,
\beta+n
).
}
\]

For a normal likelihood with known variance and normal prior,

\[
\boxed{
\frac1{\tau_n^2}
=
\frac1{\tau_0^2}
+
\frac{n}{\sigma^2}.
}
\]

Posterior prediction uses

\[
\boxed{
p
(
x_{\text{new}}
\mid x
)
=
\int
p
(
x_{\text{new}}
\mid\theta
)
\pi(\theta\mid x)
\,d\theta.
}
\]

Bayesian decision theory chooses actions by minimizing

\[
\boxed{
\mathbb{E}
[
L(\theta,a)
\mid x
].
}
\]

Posterior odds satisfy

\[
\boxed{
\text{posterior odds}
=
\text{prior odds}
\times
\text{Bayes factor}.
}
\]

Hierarchical models create partial pooling and shrinkage across related
groups.

Modern Bayesian computation uses:

\[
\boxed{
\text{Monte Carlo}
+
\text{MCMC}
+
\text{optimization}
+
\text{approximation}.
}
\]

Thus Bayesian statistics combines

\[
\boxed{
\text{probability}
+
\text{likelihood}
+
\text{prior information}
+
\text{decision theory}
+
\text{computation}.
}
\]

%%%%%%%%%%%%%%%%%%%%%%%%%%%%%%%%%%%%%%%%%%%%%%%%%%%%%%%%%%%%
\subsection{Section Connections}
\label{subsec:bayesian-connections}
%%%%%%%%%%%%%%%%%%%%%%%%%%%%%%%%%%%%%%%%%%%%%%%%%%%%%%%%%%%%

\chapterconnection{
Bayesian Statistics builds directly on the likelihood, decision-theoretic,
and asymptotic foundations developed in
Section~\ref{sec:mathematical-statistics}. Bayes' theorem transforms
prior distributions and likelihoods into posterior distributions,
providing probability statements about parameters and future
observations. Hierarchical models extend regression, ANOVA, and
multivariate models through partial pooling, while prior distributions
connect naturally with regularization methods from statistical learning.
Posterior computation introduces Monte Carlo, Markov chain Monte Carlo,
and optimization methods that are developed more fully in Computational
Statistics. The next section, Time Series Analysis, studies data indexed
through time and introduces statistical dependence across ordered
observations.
}

%%%%%%%%%%%%%%%%%%%%%%%%%%%%%%%%%%%%%%%%%%%%%%%%%%%%%%%%%%%%
% SECTION SOURCE TRACKING
%%%%%%%%%%%%%%%%%%%%%%%%%%%%%%%%%%%%%%%%%%%%%%%%%%%%%%%%%%%%

%------------------------------------------------------------
% 8.14 Bayesian Statistics
%------------------------------------------------------------
%
% Source basis:
%   - Standard Bayesian statistical theory.
%   - Standard conjugate Bayesian models.
%   - Standard Bayesian decision theory.
%   - Standard Bayesian computational concepts.
%
% Used for:
%   - prior distributions;
%   - likelihoods;
%   - posterior distributions;
%   - marginal likelihoods;
%   - discrete Bayesian updating;
%   - conjugate priors;
%   - Beta--Bernoulli models;
%   - Beta--Binomial prediction;
%   - Gamma--Poisson models;
%   - Normal--Normal models;
%   - posterior means, medians, and modes;
%   - MAP estimation;
%   - credible intervals;
%   - highest posterior density regions;
%   - posterior predictive distributions;
%   - informative and weakly informative priors;
%   - diffuse and improper priors;
%   - posterior propriety;
%   - prior predictive checks;
%   - posterior predictive checks;
%   - Bayesian model checking;
%   - posterior odds;
%   - Bayes factors;
%   - Bayesian decision theory;
%   - loss functions and Bayes actions;
%   - hierarchical models;
%   - hyperparameters;
%   - complete, no, and partial pooling;
%   - exchangeability;
%   - empirical Bayes;
%   - Bayesian regression;
%   - Bayesian regularization;
%   - Bayesian logistic regression;
%   - Monte Carlo integration;
%   - Markov chain Monte Carlo;
%   - Metropolis--Hastings;
%   - Gibbs sampling;
%   - Hamiltonian Monte Carlo preview;
%   - NUTS preview;
%   - MCMC autocorrelation;
%   - effective sample size;
%   - Monte Carlo standard error;
%   - trace plots;
%   - convergence diagnostics;
%   - importance sampling preview;
%   - Laplace approximation;
%   - variational inference;
%   - KL divergence;
%   - ELBO;
%   - prior sensitivity;
%   - Bayesian hypothesis testing;
%   - point null hypotheses;
%   - practical equivalence regions;
%   - Bayesian model averaging;
%   - applications and exercises.
%
% Notes:
%   - Time Series Analysis is developed next in Section 8.15.
%   - Computational Statistics later develops MCMC, importance sampling,
%     bootstrap, optimization, and simulation methods in greater depth.
%   - Statistical Learning develops Bayesian regularization and
%     predictive modeling more deeply.
%
%%%%%%%%%%%%%%%%%%%%%%%%%%%%%%%%%%%%%%%%%%%%%%%%%%%%%%%%%%%%